\documentclass[11pt,a4paper]{article}

\usepackage[utf8]{inputenc}
\usepackage[T1]{fontenc}
\usepackage{lmodern}
\usepackage{microtype}
\usepackage[a4paper,margin=1in]{geometry}
\usepackage{amsmath,amssymb,amsthm,mathtools}
\usepackage{mathrsfs}
\usepackage{booktabs}
\usepackage{enumitem}
\usepackage{array}
\usepackage{xcolor}
\usepackage{orcidlink}
\usepackage{hyperref}
\usepackage[nameinlink,noabbrev]{cleveref}

\hypersetup{
  colorlinks=true,
  linkcolor=blue!55!black,
  citecolor=blue!55!black,
  urlcolor=blue!55!black,
  pdftitle={Homothetic Valuation Spectrum of a Quasilocal B4 Filling and an Effective Residual Sector on Compact S3},
  pdfauthor={Andrei Preece and Boris Batenin},
  pdfsubject={Closed cosmology; quasilocal filling; effective residual sector},
  pdfkeywords={closed cosmology; S3; B4 filling; quasilocal energy; residual dark sector; valuation spectrum; Noether identities; S4 compensation}
}

% --- Fix for hyperref warnings: safely expand math macros in PDF strings ---
\pdfstringdefDisableCommands{%
  \def\Bfour{B4}%
  \def\Sthree{S3}%
  \def\Sfour{S4}%
  \def\dd{d}%
}

\setlist[itemize]{leftmargin=*,itemsep=0.2em,topsep=0.4em}
\setlist[enumerate]{leftmargin=*,itemsep=0.25em,topsep=0.4em}
\renewcommand{\arraystretch}{1.15}
\numberwithin{equation}{section}

\newtheorem{criterion}{Criterion}
\newtheorem{definition}[criterion]{Definition}
\newtheorem{lemma}[criterion]{Lemma}
\newtheorem{theorem}[criterion]{Theorem}
\newtheorem{proposition}[criterion]{Proposition}
\newtheorem{corollary}[criterion]{Corollary}
\newtheorem{hypothesis}[criterion]{Hypothesis}
\newtheorem{closurecondition}[criterion]{Closure condition}
\newtheorem{nogotheorem}[criterion]{No-go theorem}
\theoremstyle{remark}
\newtheorem{remark}[criterion]{Remark}

\newcommand{\Sthree}{S^3}
\newcommand{\Sfour}{S^4}
\newcommand{\Bfour}{B^4}
\newcommand{\M}{\mathcal{M}}
\newcommand{\Dom}{\mathcal{U}}
\newcommand{\Dil}{\mathscr{D}}
\newcommand{\B}{\mathcal{B}}
\newcommand{\Efour}{E_4}
\newcommand{\rhoFour}{\rho_4}
\newcommand{\epsD}{\varepsilon_D}
\newcommand{\pD}{p_D}
\newcommand{\QD}{Q_D}
\newcommand{\PiA}{\Pi_A}
\newcommand{\vol}{\operatorname{Vol}}
\newcommand{\diag}{\operatorname{diag}}
\newcommand{\dd}{\mathrm{d}}

\title{Homothetic Valuation Spectrum of a Quasilocal \(\Bfour\) Filling\\
       and an Effective Residual Sector on Compact \(\Sthree\)}
\author{Andrei~Preece\,\orcidlink{0009-0004-5978-795X}\\
        \small \href{mailto:andre.preece@web.de}{andre.preece@web.de}\\
        \small Independent Researcher, Berlin, Germany
        \and
        Boris~Batenin\,\orcidlink{0009-0000-7027-8583}\\
        \small \href{mailto:batenin@web.de}{batenin@web.de}\\
        \small Independent Researcher, Berlin, Germany}
\date{2026}

\begin{document}
\maketitle


\begin{abstract}
We develop a geometric-variational representation of homogeneous residual sectors on a closed FLRW spacetime \(I\times \Sthree_A\), using the Riemannian ball \(\Bfour_A\) only as an auxiliary quasilocal filling rather than as a physical extra dimension. For the round homothetic filling,
\begin{equation*}
V_3=2\pi^2A^3,
\qquad
V_4=\frac{\pi^2}{2}A^4,
\end{equation*}
a filling energy \(\Efour=\rhoFour V_4\) induces the boundary density and adiabatic pressure
\begin{equation*}
\epsD=\frac{\Efour}{V_3}=\frac A4\rhoFour,
\qquad
\pD=-\frac{\dd\Efour}{\dd V_3}.
\end{equation*}
The radial derivative \(V_3^{-1}\dd\Efour/\dd A\) is a conjugate response, not a second boundary density. The construction yields the conservation identity, the inverse representation \(\rhoFour=4\epsD/A\), the local equation-of-state relation
\begin{equation*}
w_D(A)=-1-\frac13\frac{\dd\ln\epsD}{\dd\ln A},
\end{equation*}
and the dimensional scaling law \(w_d=(\beta-d-1)/d\) for power-law \((d+1)\)-dimensional fillings.

Under Euclidean-invariant homothetic valuation assumptions for the auxiliary convex filling in \(\mathbb{R}^4\), the full valuation ansatz contains a bulk intrinsic-volume term and four boundary curvature terms,
\begin{equation*}
\Efour^{\rm val}(A)=\gamma V_4(A)+\sum_{j=0}^{3}\alpha_j
\int_{\Sthree_A}e_j(K)\,\dd\Sigma .
\end{equation*}
For the round representative this gives
\begin{equation*}
\epsD(A)=
\frac{\gamma A}{4}
+\alpha_0+\frac{3\alpha_1}{A}
+\frac{3\alpha_2}{A^2}
+\frac{\alpha_3}{A^3},
\end{equation*}
with branch equations of state
\begin{equation*}
w=-\frac43,\ -1,\ -\frac23,\ -\frac13,\ 0.
\end{equation*}
Equivalently, these are the five eigenbranches \(s=1,0,-1,-2,-3\) of the homothetic dilation operator \(\Dil=A\,\dd/\dd A\). Membership in the round homothetic profile image of the full valuation class is characterized sharply by
\begin{equation*}
(\Dil-1)\Dil(\Dil+1)(\Dil+2)(\Dil+3)\epsD=0,
\end{equation*}
while the minimal endpoint density profile is characterized by \(\Dil(\Dil+3)\epsD=0\); the lift to an endpoint functional requires the separate minimal endpoint assumptions.
A separate minimal endpoint classification is imposed on local \({\rm Diff}(\Sthree)\)-invariant boundary functionals with no local endpoint degrees of freedom, no curvature duplication, no fixed external structures, additivity, and internal-field dependence only through an oriented topological charge. Inside this declared class,
\begin{equation*}
E_{\rm end}(A,Q)=\sigma V_3(A)+\mu Q,
\end{equation*}
which induces
\begin{equation*}
\epsD(a)=\varepsilon_q+\varepsilon_w a^{-3},
\qquad
\pD=-\varepsilon_q .
\end{equation*}
The endpoint form is therefore not a theorem from topology alone; it is conditionally unique only inside the stated minimal endpoint class. A minimal four-dimensional effective completion realizes this sector as vacuum energy plus a conserved pressureless current, with \(c_s^2=0\), \(\pi_{\mu\nu}=0\), and geodesic flow. The result is a geometric residual classification and minimal effective completion framework, not a topology-only derivation of the observed dark-sector amplitudes. The construction classifies and selects functional forms, but it does not determine the dimensional amplitudes, a microscopic matter model, transfer functions, likelihoods, or a complete dynamics for \(A(\tau)\).
\end{abstract}

\noindent\textbf{Keywords:} closed cosmology; \(\Sthree\); \(\Bfour\) filling; quasilocal energy; effective residual dark sector; Noether identities; \(\Sfour\) compensation.

\section*{Summary of the Result}
\begin{center}
\begin{tabular}{@{}p{0.27\linewidth}p{0.66\linewidth}@{}}
\toprule
Established geometry & Round \(\Sthree_A\), homothetic \(\Bfour_A\), quasilocal cones, \(V_4/V_3=A/4\), and orientation-sensitive \(\Sfour\) gluing. \\
Derived background results & Direct and inverse reconstruction, pressure from work, conservation, local \(w_D(A)\), dimensional scaling, and linear superposition. \\
Full valuation spectrum & The homothetic valuation class gives the five-slot spectrum \(w=-4/3,-1,-2/3,-1/3,0\). \\
Dilation-spectrum closure & The round profile image is exactly \(\ker[(\Dil-1)\Dil(\Dil+1)(\Dil+2)(\Dil+3)]\); spectral projectors extract the five invariant branch amplitudes. \\
Sharp endpoint test & The minimal endpoint density is exactly \(\ker[\Dil(\Dil+3)]\), with surviving amplitudes \(\varepsilon_q=\mathcal I_0\) and \(\varepsilon_w=\mathcal I_{-3}\). \\
Pressure kernel & \(p=0\) iff the endpoint energy is constant in \(A\); division by \(V_3\) maps this kernel to the dust scaling \(a^{-3}\). \\
Boundary-only valuation filter & Within the boundary-curvature valuation subclass, \(\gamma=0\) removes the bulk \(V_4\) branch. \\
Endpoint selection & Inside the stated minimal endpoint class, \(E_{\rm end}=\sigma V_3+\mu Q\), hence \(\epsD=\varepsilon_q+\varepsilon_w a^{-3}\). \\
Model conditions & Matching \(T^{\rm fill}_{\mu\nu}=T^{\rm rec}_{\mu\nu}\), consistent exchange, minimal endpoint assumptions, perturbative completion, and independent amplitude input or calibration. \\
Protected branches & \(\Sfour\) compensation, Hopf modes, causal accessibility, radius protection, and non-minimal perturbative completions. \\
Completion level & A minimal vacuum plus conserved-dust effective completion exists; conditional endpoint uniqueness, amplitudes, the scope of the minimal endpoint class, and observational inference remain separate. \\
\bottomrule
\end{tabular}
\end{center}

\tableofcontents

\section{Position of the Work and Scope of the Result}
This paper is a local mathematical node of a broader research corpus. The wider ontological program motivates the construction but is not a theorem-level output of the present work.

The proved core of the present work is the chain
\begin{equation}
\Sthree_A\longrightarrow \Bfour_A\ \text{filling}\longrightarrow \Efour(A,\xi)\longrightarrow (\epsD,\pD,\QD).
\end{equation}
The \(\Sfour\)-compensation construction is retained as a downstream branch. Only orientational cancellation under compatible gluing is proved here; a dynamical compensation law and the selection of physical amplitudes remain open.

This chain is not presented as an autonomous brane model, a standard extra-dimensional cosmology, or a physical embedding of the observed space into a five-dimensional bulk. Similarity between an isolated formula and a formula in another model does not establish structural equivalence. The adjective quasilocal is used here only in the sense that a bounded spatial domain is assigned an auxiliary filling energy. It does not imply the Brown-York boundary stress tensor, the Hamilton-Jacobi gravitational energy, or a physical bulk embedding. In particular, \(\Efour\) is not identified with the Brown-York gravitational quasilocal energy \cite{BrownYork1993}, whose variational origin and boundary stress tensor are structurally different. No Hamiltonian boundary stress tensor is constructed from \(\Efour\); the term quasilocal refers only to the assignment of an auxiliary filling energy to a spatial domain.

The precise result has two directions. The direct construction produces an effective background fluid from a filling energy. The inverse representation theorem shows that every differentiable homogeneous background density has a unique representing \(\Bfour\) density within the adopted homothetic class; a prescribed pair \((\epsD,\pD)\) is represented by the adiabatic map exactly when it satisfies the corresponding continuity identity. This establishes both the strength of the construction and the exact limit of the universal map: geometry fixes the representation. Predictivity begins only after an invariant closure or additional dynamical principle selects a restricted functional class for \(\Efour\).

\paragraph{Statement of results.}
The central theorem-level content is the finite spectral closure of the round homothetic profile image of the valuation class under the dilation operator \(\Dil=A\,\dd/\dd A\). Its round homothetic profile image is characterized exactly by
\begin{equation}
(\Dil-1)\Dil(\Dil+1)(\Dil+2)(\Dil+3)\epsD=0,
\end{equation}
so the five background branches are the five dilation eigenvalues \(s=1,0,-1,-2,-3\). The boundary-only and minimal endpoint subclasses then admit sharp obstruction tests, with the minimal endpoint density subspace given by
\begin{equation}
\Dil(\Dil+3)\epsD=0
\quad\Longleftrightarrow\quad
\epsD=\varepsilon_q+\varepsilon_w a^{-3}.
\end{equation}
Inside the explicitly stated minimal endpoint class, this density is realized by \(E_{\rm end}=\sigma V_3+\mu Q\). The selected sector admits a minimal four-dimensional effective completion as curved \(\Lambda\)CDM with cold dust perturbations. Observational superiority is not assumed; it requires background integrability, independent compact-\(\Sthree\) spectral-geometric support, and robust penalized statistical performance against the nearest competing cosmologies.

\begin{criterion}[Minimal physical completeness]
A candidate physical dark sector must provide more than a background function \(H(a)\). It must also specify an energy-momentum tensor, a conservation or exchange law, and rules for perturbations. The background map is a necessary but not sufficient level of a physical model.
\end{criterion}

\section{Conventions and Geometric Arena}
\subsection{Carrier, Metric Representative, and Reconstruction Level}
\begin{definition}[Carrier, metric representative, and reconstruction level]
Let
\begin{equation}
\Sigma_{\rm car}\simeq \Sthree
\end{equation}
denote the underlying topological carrier. For each \(A>0\),
\begin{equation}
(\Sthree,g_A),\qquad g_A=A^2\dd\Omega_3^2,
\end{equation}
is its round homothetic metric representative. For an observer \(O\) at proper time \(\tau\),
\begin{equation}
\Dom_O(\tau)\subseteq \Sthree_A(\tau)
\end{equation}
denotes the causally accessible or operationally reconstructed domain.
\end{definition}

These three levels are not identified. In particular, \(A(\tau)\) is a scale parameter of the metric representative. The present background map does not by itself interpret limiting values of \(A\) as creation, destruction, collapse, or rupture of \(\Sigma_{\rm car}\).

We use metric signature \((-+++)\). The physical arena is the closed \(k=+1\) FLRW class,
\begin{equation}
\M=I\times \Sigma,
\qquad
\Sigma\simeq \Sthree,
\qquad
A(t)=a(t)R_0.
\end{equation}
We use the normalization \(a(t_0)=1\), so that
\begin{equation}
A_0=A(t_0)=R_0.
\end{equation}
The metric representative is
\begin{equation}
\dd s^2=-N(t)^2c^2\dd t^2+A(t)^2\dd\Omega_3^2.
\end{equation}
This is the standard closed FLRW convention used in relativistic cosmology \cite{HawkingEllis1973,MTW1973,Weinberg2008}.
Proper time satisfies \(\dd\tau=N\dd t\), and
\begin{equation}
H=\frac{1}{A}\frac{\dd A}{\dd\tau}=\frac{\dot A}{NA},
\end{equation}
where the dot denotes differentiation with respect to coordinate time \(t\). Derivatives with respect to proper time are written explicitly as \(\dd/\dd\tau\).

For the round three-sphere,
\begin{equation}
V_3(A)=\vol(\Sthree_A)=2\pi^2A^3,
\qquad
{}^{(3)}R=\frac{6}{A^2}.
\end{equation}
The Riemannian filling \(\Bfour_A\) has
\begin{equation}
\partial \Bfour_A=\Sthree_A,
\qquad
V_4(A)=\vol(\Bfour_A)=\frac{\pi^2}{2}A^4,
\qquad
\frac{\dd V_4}{\dd A}=V_3.
\end{equation}
Hence
\begin{equation}
\frac{V_4}{V_3}=\frac{A}{4}.
\end{equation}
The factor \(A/4\) is a metric consequence of the round homothetic filling, not a consequence of the topological relation \(\Sthree=\partial\Bfour\) alone.

\begin{remark}[Auxiliary status of the filling]
The Riemannian ball \(\Bfour_A\) is attached to the metric representative \((\Sthree,g_A)\), not to the bare topological carrier alone. It is an auxiliary filling used to define a geometric measure and a variational energy; it is neither an additional physical spatial dimension nor a physical five-dimensional spacetime bulk.
\end{remark}

The \(\Bfour\)-filling can geometrically represent any separately conserved homogeneous background component within the adopted homothetic class, but it does not determine its microscopic origin, perturbations, sound speed, anisotropic stress, particle content, or full spacetime dynamics.

\begin{criterion}[Compactness stress test]
Every residual-sector hypothesis should be testable both in the global compact class \(\Sthree_A\) and in the local flat limit \(A\to\infty\), \(\Sthree_A\to\mathbb{R}^3\) at fixed finite observational scale, understood as a local or pointed geometric limit rather than a change of global topology. A construction that requires openness must identify the physical term or boundary condition that requires it.
\end{criterion}

\section{Quasilocal Map and Geodesic Domains}
For a domain \(\Dom\subset \Sthree_A\), consider the quasilocal map
\begin{equation}
\Dom\longmapsto \B(\Dom)\subset \Bfour_A.
\end{equation}
Let \(\Dom_\chi\) be a geodesic ball of angular radius \(0\le \chi\le\pi\). Then
\begin{align}
V_3(\chi,A)&=4\pi A^3\int_0^\chi \sin^2u\,\dd u
=2\pi A^3(\chi-\sin\chi\cos\chi),\\
A_2(\chi,A)&=4\pi A^2\sin^2\chi.
\end{align}
Define the homothetic cone by transporting the same angular domain homothetically to each intermediate radius:
\begin{equation}
\B(\Dom_\chi):=\{(r,\Omega)\mid 0\le r\le A,\ \Omega\in \Dom_\chi^{\rm ang}\},
\end{equation}
where \(\Dom_\chi^{\rm ang}\subset \Sthree\) denotes the angular geodesic ball of radius \(\chi\) on the unit sphere.

\begin{lemma}[Volume of the conical filling]
For fixed \(\chi\),
\begin{equation}
V_4(\chi,A)=\frac{A}{4}V_3(\chi,A),
\qquad
\frac{\partial V_4(\chi,A)}{\partial A}=V_3(\chi,A).
\end{equation}
\end{lemma}
\begin{proof}
Write \(V_3(\chi,r)=C(\chi)r^3\), where \(C(\chi)=2\pi(\chi-\sin\chi\cos\chi)\). Then
\begin{equation}
V_4(\chi,A)=\int_0^A V_3(\chi,r)\,\dd r
=C(\chi)\int_0^A r^3\,\dd r
=\frac{C(\chi)A^4}{4}
=\frac{A}{4}V_3(\chi,A).
\end{equation}
Differentiation with respect to \(A\) gives the second identity.
\end{proof}

The full boundary of \(\B(\Dom_\chi)\) contains the outer domain \(\Dom_\chi\), the lateral conical part over \(\partial\Dom_\chi\), and their codimension-two corner. The conical-volume identity does not remove possible lateral or corner contributions in an action variation.

\section{Filling Energy and Two Distinct Responses}
Let \(\rhoFour(A)\) be the auxiliary filling density per spatial four-volume, \([\rhoFour]=EL^{-4}\). Define
\begin{equation}
\Efour(A):=\rhoFour(A)V_4(A).
\end{equation}

\begin{definition}[Accumulated boundary density]
\begin{equation}
\epsD(A):=\frac{\Efour(A)}{V_3(A)}=\frac{A}{4}\rhoFour(A).
\end{equation}
\end{definition}

\begin{definition}[Radially conjugate response]
\begin{equation}
\PiA(A):=\frac{1}{V_3(A)}\frac{\dd\Efour}{\dd A}
=\rhoFour(A)+\frac{A}{4}\rhoFour'(A).
\end{equation}
\end{definition}

The dimensions are different:
\begin{equation}
[\epsD]=EL^{-3},
\qquad
[\PiA]=EL^{-4}.
\end{equation}
Thus \(\PiA\) is not a second three-dimensional energy density and is not the physical pressure. It is a radial force density, or equivalently an energetic response per unit boundary three-volume. The physical three-dimensional thermodynamic pressure conjugate to \(V_3\) is instead \(\pD=-(A/3)\PiA\).

For homogeneous \(\rhoFour(A)\), the quasilocal density on every fixed \(\Dom_\chi\) agrees with the global one:
\begin{equation}
\frac{\Efour(\Dom_\chi,A)}{V_3(\chi,A)}
=\frac{\rhoFour(A)V_4(\chi,A)}{V_3(\chi,A)}
=\frac{A}{4}\rhoFour(A).
\end{equation}
The independence from \(\chi\) follows here from homogeneity and homothety; it is not a general law for arbitrary fillings.

\section{Pressure and Direct Reconstruction}
\begin{definition}[Adiabatic pressure]
For a one-parameter energy \(\Efour=\Efour(A)\),
\begin{equation}
\pD(A):=-\frac{\dd\Efour}{\dd V_3}
=-\frac{\dd\Efour/\dd A}{\dd V_3/\dd A}.
\end{equation}
Since \(\dd V_3/\dd A=3V_3/A\),
\begin{equation}
\pD=-\frac{A}{3}\PiA
=-\frac{A}{3}\left(\rhoFour+\frac{A}{4}\rhoFour'\right).
\end{equation}
The factor \(A/3\) is essential: \(\PiA\) is the radial conjugate response per boundary three-volume, whereas \(\pD\) is the thermodynamic pressure conjugate to the physical boundary volume \(V_3\).
\end{definition}

\begin{theorem}[Direct background reconstruction]
Let \(A>0\), let \(\Efour(A)\) be a \(C^1\) function on this domain, and assume the round volume law \(V_3(A)=2\pi^2A^3\), so that \(V_3'(A)=3V_3/A\neq0\). Then \(\epsD=\Efour/V_3\) and \(\pD=-\dd\Efour/\dd V_3\) uniquely define an isotropic background fluid satisfying
\begin{equation}
\frac{\dd\epsD}{\dd\tau}+3H(\epsD+\pD)=0.
\end{equation}
\end{theorem}
\begin{proof}
Write \(V=V_3(A)\) and \(E=\Efour(A)\). Since \(V'/V=3/A\), one has
\begin{equation}
\frac{\dd\epsD}{\dd A}
=\frac{E'}{V}-\frac{EV'}{V^2}
=\frac{E'}{V}-\frac{3E}{AV}.
\end{equation}
Moreover,
\begin{equation}
\pD=-\frac{E'}{V'}
=-\frac{AE'}{3V},
\end{equation}
and therefore
\begin{equation}
\frac{3}{A}(\epsD+\pD)
=\frac{3E}{AV}-\frac{E'}{V}.
\end{equation}
Adding the last two equations gives
\begin{equation}
\frac{\dd\epsD}{\dd A}
+\frac{3}{A}(\epsD+\pD)=0.
\end{equation}
Multiplication by \(\dd A/\dd\tau=HA\) gives the continuity identity. Uniqueness follows from the definitions.
\end{proof}

\section{Internal Variables and Exchange}
Let
\begin{equation}
\Efour=\Efour(A,\xi^1,\ldots,\xi^n).
\end{equation}
Pressure is defined at fixed internal variables:
\begin{equation}
\pD:=-\left.\frac{\partial \Efour}{\partial V_3}\right|_{\xi}.
\end{equation}

\begin{proposition}[Exchange term]
Under evolution \(\xi^I(\tau)\),
\begin{equation}
\frac{\dd\epsD}{\dd\tau}+3H(\epsD+\pD)=\QD,
\qquad
\QD=\frac{1}{V_3}\sum_I\frac{\partial\Efour}{\partial\xi^I}\frac{\dd\xi^I}{\dd\tau}.
\end{equation}
\end{proposition}
\begin{proof}
The total derivative is
\begin{equation}
\frac{\dd\Efour}{\dd\tau}
=\frac{\partial\Efour}{\partial A}\frac{\dd A}{\dd\tau}
+\sum_I\frac{\partial\Efour}{\partial\xi^I}\frac{\dd\xi^I}{\dd\tau}.
\end{equation}
The first terms cancel exactly as in the proof of the continuity identity; the remainder is the exchange law.
\end{proof}

Thus \(\QD\) is not arbitrary within this class: it measures the nonadiabatic work of the internal filling variables.

\section{Inverse Homogeneous Representation and the Limit of Background Predictivity}
\begin{theorem}[Inverse background representation]
Let \(\epsD(A)\) be a differentiable background density. Within the homogeneous homothetic filling class there is a unique representing density
\begin{equation}
\rhoFour(A)=\frac{4}{A}\epsD(A),
\end{equation}
such that \(\Efour=\rhoFour V_4=\epsD V_3\). The pressure induced by the filling-work prescription is
\begin{equation}
\pD(A)=-\epsD(A)-\frac{A}{3}\frac{\dd\epsD}{\dd A}.
\end{equation}
Conversely, a prescribed pair \((\epsD,\pD)\) is representable by this adiabatic filling map if and only if it satisfies
\begin{equation}
\pD(A)=-\epsD(A)-\frac{A}{3}\epsD'(A),
\end{equation}
equivalently the homogeneous continuity identity.
\end{theorem}
\begin{proof}
The relation \(\Efour=\epsD V_3\) together with \(V_3/V_4=4/A\) immediately gives the inverse relation. The same volume ratio gives uniqueness. Furthermore,
\begin{equation}
\pD=-\frac{\dd(\epsD V_3)}{\dd V_3}
=-\epsD-V_3\frac{\dd\epsD/\dd A}{\dd V_3/\dd A}
=-\epsD-\frac{A}{3}\epsD',
\end{equation}
which is the induced pressure relation. The final equivalence follows by multiplying \(\epsD'+3(\epsD+\pD)/A=0\) by \(\dd A/\dd\tau\).
\end{proof}

\begin{remark}[Exact scope of inverse representation]
The inverse representation map
\begin{equation}
\epsD(A)\longmapsto \rhoFour(A)=\frac{4}{A}\epsD(A)
\end{equation}
is homogeneous and algebraic. It establishes uniqueness of the representing density only within the adopted homothetic background class and on the \(A\)-domain where the background history is expressed as a differentiable function of \(A\). It does not reconstruct microscopic fields, particles, perturbations, or a covariant matter action.
\end{remark}

\begin{remark}[No inverse reconstruction of perturbations]
The inverse map reconstructs only the homogeneous representative \(\rhoFour(A)=4\epsD(A)/A\). It is not an inverse map from a background history to a unique stress-energy tensor, perturbation theory, sound speed, anisotropic stress, or microscopic action. These require a separate covariant completion.
\end{remark}

Any interpretation of a singular limit as a boundary of inverse physical reconstruction is an additional interpretive principle, not a consequence of the inverse background representation theorem proved here. Predictivity begins when \(\Efour(A,\xi)\) is independently fixed by a symmetry, an action, a compensation condition, a spectral principle, or equations for the internal variables.

\section{General Equation-of-State Formula}
For \(\epsD>0\), define \(w_D=\pD/\epsD\). The inverse pressure relation gives
\begin{equation}
w_D(A)=-1-\frac{1}{3}\frac{\dd\ln\epsD}{\dd\ln A}
=-\frac{4}{3}-\frac{1}{3}\frac{\dd\ln\rhoFour}{\dd\ln A}.
\end{equation}
Define the local scaling exponent
\begin{equation}
\beta_{\rm eff}(A):=-\frac{\dd\ln\rhoFour}{\dd\ln A}.
\end{equation}
Then
\begin{equation}
w_D(A)=\frac{\beta_{\rm eff}(A)-4}{3}.
\end{equation}

\begin{corollary}[Power-law class]
If
\begin{equation}
\rhoFour(A)=\rho_{4,0}\left(\frac{A}{A_0}\right)^{-\beta},
\end{equation}
then
\begin{equation}
\epsD\propto A^{1-\beta},
\qquad
\pD=\frac{\beta-4}{3}\epsD,
\qquad
w_D=\frac{\beta-4}{3}.
\end{equation}
\end{corollary}

\begin{center}
\begin{tabular}{@{}cccc@{}}
\toprule
\(\beta\) & Scaling of \(\epsD\) & \(w_D\) & Background class \\
\midrule
0 & \(a\)      & \(-4/3\) & phantom-like slot \\
1 & \(a^0\)    & \(-1\)   & vacuum slot \\
2 & \(a^{-1}\) & \(-2/3\) & accelerating slot \\
3 & \(a^{-2}\) & \(-1/3\) & boundary slot \\
4 & \(a^{-3}\) & \(0\)    & dust slot \\
5 & \(a^{-4}\) & \(1/3\)  & radiation-like slot \\
\bottomrule
\end{tabular}
\end{center}

For \(\epsD>0\),
\begin{equation}
\epsD+
\pD=\frac{\beta-1}{3}\epsD,
\qquad
\epsD+3\pD=(\beta-3)\epsD.
\end{equation}
Thus \(\beta<1\) violates the null energy condition of the effective fluid, while \(\beta<3\) gives a negative active gravitational density \(\epsD+3\pD\). These background properties do not replace a microscopic stability analysis.

\section{Dimensional Theorem}
\begin{theorem}[Homothetic filling in arbitrary dimension]
Let
\begin{equation}
V_d(A)=C_dA^d,
\qquad
V_{d+1}(A)=C_{d+1}A^{d+1},
\qquad
\frac{\dd V_{d+1}}{\dd A}=V_d,
\end{equation}
so that \(C_d=(d+1)C_{d+1}\). Define
\begin{equation}
E_{d+1}=\rho_{d+1}V_{d+1},
\qquad
\varepsilon_d=\frac{E_{d+1}}{V_d},
\qquad
p_d=-\frac{\dd E_{d+1}}{\dd V_d}.
\end{equation}
Then
\begin{equation}
\varepsilon_d=\frac{A}{d+1}\rho_{d+1},
\qquad
p_d=-\frac{A}{d}\left(\rho_{d+1}+\frac{A}{d+1}\rho_{d+1}'\right).
\end{equation}
If \(\rho_{d+1}\propto A^{-\beta}\), then
\begin{equation}
w_d=\frac{\beta-d-1}{d}.
\end{equation}
\end{theorem}
\begin{proof}
The identity \(\dd V_{d+1}/\dd A=V_d\) gives \(V_{d+1}/V_d=A/(d+1)\), which proves the dimensional density relation. Since \(V_d'/V_d=d/A\),
\begin{equation}
p_d=-\frac{A}{dV_d}\frac{\dd E_{d+1}}{\dd A}.
\end{equation}
Differentiating \(E_{d+1}=\rho_{d+1}C_{d+1}A^{d+1}\) gives the dimensional pressure relation. For a power-law density, \(\rho_{d+1}'=-(\beta/A)\rho_{d+1}\); division by the dimensional density relation yields the dimensional equation-of-state relation.
\end{proof}

The coefficients \(1/4\) and \(4\) in the physical case \(d=3\) are therefore special values of a general dimensional structure.

\section{Linear Superposition and Protected Slots}
\begin{theorem}[Linear superposition]
If \(\Efour(A)=\sum_i E_{4,i}(A)\), then
\begin{equation}
\epsD=\sum_i \varepsilon_{D,i},
\qquad
\pD=\sum_i p_{D,i}.
\end{equation}
\end{theorem}
\begin{proof}
Both maps, \(\Efour\mapsto \Efour/V_3\) and \(\Efour\mapsto -\dd\Efour/\dd V_3\), are linear.
\end{proof}

For
\begin{equation}
\rhoFour(A)=\sum_i C_iA^{-\beta_i},
\end{equation}
one obtains
\begin{equation}
\epsD(A)=\sum_i \frac{C_i}{4}A^{1-\beta_i},
\qquad
\pD(A)=\sum_i \frac{\beta_i-4}{3}\varepsilon_{D,i}(A).
\end{equation}
The two-branch law
\begin{equation}
\rhoFour(A)=\frac{4\varepsilon_q}{A}+\frac{4\varepsilon_wA_0^3}{A^4}
\end{equation}
is especially important. With \(a=A/A_0\),
\begin{equation}
\epsD(a)=\varepsilon_q+\varepsilon_w a^{-3},
\qquad
\pD(a)=-\varepsilon_q.
\end{equation}
Define the diagnostic invariants
\begin{equation}
\mathcal{A}_q(a):=-\pD(a),
\qquad
\mathcal{A}_w(a):=a^3\bigl(\epsD(a)+\pD(a)\bigr).
\end{equation}

\begin{criterion}[Exact vacuum-dust decomposition]
At the background level, the conditions
\begin{equation}
\frac{\dd\mathcal{A}_q}{\dd a}=0,
\qquad
\frac{\dd\mathcal{A}_w}{\dd a}=0
\end{equation}
are necessary and sufficient for the representation with constant \(\varepsilon_q\) and \(\varepsilon_w\).
\end{criterion}
\begin{proof}
Necessity follows by direct substitution. Conversely, constancy of \(\mathcal{A}_q\) gives \(\pD=-\varepsilon_q\), while constancy of \(\mathcal{A}_w\) gives \(\epsD+
\pD=\varepsilon_w a^{-3}\). Adding the two equations yields the decomposition.
\end{proof}

The dust slot is an exact background class, but it is not yet clustering cold dark matter. That identification requires perturbative closure.



\section{Full Valuation Spectrum, Boundary Restriction, and Minimal Endpoint Classification}
The universal background map
\begin{equation}
\Efour(A)\longmapsto(\epsD,\pD)
\end{equation}
is algebraically complete but not predictive by itself. Predictivity begins only after a restricted invariant class of filling energies has been selected.

We do not reprove Hadwiger's theorem here \cite{Hadwiger1957,Schneider2014}. We use its intrinsic-volume classification only to motivate the finite valuation basis and then compute its round homothetic specialization explicitly.

\begin{hypothesis}[Euclidean convex valuation closure]
The auxiliary filling energy is assumed to be a continuous, Euclidean-invariant, additive valuation on convex bodies in \(\mathbb{R}^4\), restricted to the round homothetic family \(\Bfour_A\). For the smooth round representatives considered here, this intrinsic-volume span may be written, after absorbing conventional normalizations into the coefficients, as
\begin{equation}
\Efour^{\rm val}(A)
=\gamma V_4(A)
+\sum_{j=0}^{3}\alpha_j
\int_{\Sthree_A}e_j(K)\,\dd\Sigma .
\end{equation}
\end{hypothesis}

The constants \(\gamma,\alpha_j\) absorb the conventional normalizations of the intrinsic volumes and curvature measures. Only the powers of \(A\) and the induced equations of state are invariant under these normalization choices. Here \(K^a{}_b\) is the Weingarten map of the boundary hypersurface and \(e_j(K)\) are the elementary symmetric polynomials of the principal curvatures:
\begin{equation}
e_0=1,
\qquad
e_1=\operatorname{tr}K,
\qquad
e_2=\sum_{a<b}\kappa_a\kappa_b,
\qquad
e_3=\det(K^a{}_b).
\end{equation}
We use the Weingarten-map sign convention for which the outward normal of the convex round filling gives positive principal curvatures. With the opposite convention the odd elementary symmetric terms change sign; this does not change the valuation spectrum, since the coefficients \(\alpha_j\) absorb the corresponding sign convention. For the round homothetic representative,
\begin{equation}
K^a{}_b=\frac{1}{A}\delta^a_b,
\end{equation}
and hence
\begin{equation}
e_j(K)=\binom{3}{j}A^{-j}.
\end{equation}
Under the standard convex-body valuation assumptions of Euclidean invariance, continuity, and additivity, these terms correspond to the intrinsic-volume basis; the explicit formula below is the homothetic round specialization needed here.

\begin{theorem}[Full homothetic valuation spectrum]
Under the Euclidean convex valuation closure hypothesis \cite{Hadwiger1957,Schneider2014}, for the round homothetic filling \(\Bfour_A\),
\begin{equation}
\Efour^{\rm val}(A)
=\frac{\pi^2}{2}\gamma A^4
+2\pi^2\left(
\alpha_0A^3+3\alpha_1A^2+3\alpha_2A+\alpha_3
\right).
\end{equation}
The induced boundary density is
\begin{equation}
\epsD(A)
=\frac{\gamma A}{4}
+\alpha_0+\frac{3\alpha_1}{A}
+\frac{3\alpha_2}{A^2}
+\frac{\alpha_3}{A^3}.
\end{equation}
The pure valuation branches have equations of state
\begin{equation}
\gamma V_4:\quad w=-\frac{4}{3},
\end{equation}
and
\begin{equation}
e_j:\quad w_j=-1+\frac{j}{3}\qquad (j=0,1,2,3).
\end{equation}
Thus the full homothetic valuation spectrum is
\begin{equation}
w=-\frac{4}{3},\ -1,\ -\frac{2}{3},\ -\frac{1}{3},\ 0.
\end{equation}
For the general five-term valuation sector,
\begin{equation}
\pD(A)=-\frac{\gamma A}{3}-\alpha_0-\frac{2\alpha_1}{A}-\frac{\alpha_2}{A^2}.
\end{equation}
\end{theorem}

\begin{proof}
The bulk term is
\begin{equation}
\gamma V_4(A)=\frac{\pi^2}{2}\gamma A^4.
\end{equation}
For the boundary terms, since the unit three-sphere has volume \(2\pi^2\), the round boundary measure is \(\dd\Sigma=A^3\dd\Omega_3\) and
\begin{equation}
\int_{\Sthree_A}\dd\Sigma=2\pi^2A^3.
\end{equation}
Therefore
\begin{equation}
\int_{\Sthree_A}e_j(K)\,\dd\Sigma
=2\pi^2\binom{3}{j}A^{3-j}.
\end{equation}
Division by \(V_3=2\pi^2A^3\) gives the stated density. The equation-of-state formula
\begin{equation}
w=-1-\frac{1}{3}\frac{\dd\ln\epsD}{\dd\ln A}
\end{equation}
then gives \(w=-4/3\) for the bulk branch \(\epsD\propto A\), and \(w_j=-1+j/3\) for the boundary branch \(\varepsilon_{D,j}\propto A^{-j}\). Summing the branch pressures gives the stated expression for \(\pD(A)\).
\end{proof}


\subsection{Dilation-Spectrum Invariants and Sharp Closure Tests}
The finite valuation spectrum has an intrinsic differential characterization. Let
\begin{equation}
a:=\frac{A}{A_0},
\qquad
\Dil:=a\frac{\dd}{\dd a}=A\frac{\dd}{\dd A}
\end{equation}
be the homothetic dilation operator. A pure density branch
\begin{equation}
\varepsilon_s(a)=C_s a^s
\end{equation}
is an eigenfunction of \(\Dil\) with eigenvalue \(s\), and its equation of state is
\begin{equation}
w_s=-1-\frac{s}{3}.
\end{equation}
Thus the full round homothetic valuation spectrum is equivalently the finite dilation spectrum
\begin{equation}
s\in\{1,0,-1,-2,-3\}.
\end{equation}

\begin{theorem}[Dilation annihilator for the round valuation profile]
Let \(\varepsilon(a)\in C^5(J)\) on an interval \(J\subset(0,\infty)\). Then \(\varepsilon\) belongs to the round homothetic profile image of the full Euclidean valuation class if and only if
\begin{equation}
\boxed{
(\Dil-1)\Dil(\Dil+1)(\Dil+2)(\Dil+3)\varepsilon=0 .
}
\end{equation}
Equivalently,
\begin{equation}
\varepsilon(a)=
C_1a+C_0+C_{-1}a^{-1}+C_{-2}a^{-2}+C_{-3}a^{-3}.
\end{equation}
The coefficient dictionary with the valuation notation is
\begin{equation}
C_1=\frac{\gamma A_0}{4},
\qquad
C_0=\alpha_0,
\qquad
C_{-1}=\frac{3\alpha_1}{A_0},
\qquad
C_{-2}=\frac{3\alpha_2}{A_0^2},
\qquad
C_{-3}=\frac{\alpha_3}{A_0^3}.
\end{equation}
\end{theorem}
\begin{proof}
Put \(x=\ln a\). Then \(\Dil=\dd/\dd x\). The displayed equation is a constant-coefficient ordinary differential equation with simple characteristic roots
\begin{equation}
1,0,-1,-2,-3.
\end{equation}
Its solution space is therefore exactly the span of
\begin{equation}
e^x,\ 1,\ e^{-x},\ e^{-2x},\ e^{-3x},
\end{equation}
which is the stated homothetic valuation form in the variable \(a\). Conversely, every function in that span is annihilated by the polynomial in \(\Dil\).
\end{proof}

For each spectral value \(s\in\{1,0,-1,-2,-3\}\), define the spectral projector
\begin{equation}
\mathsf{P}_s
:=
\prod_{\substack{t\in\{1,0,-1,-2,-3\}\\ t\ne s}}
\frac{\Dil-t}{s-t}.
\end{equation}
On the round valuation profile,
\begin{equation}
\mathsf{P}_s\varepsilon=C_s a^s,
\qquad
\boxed{
\mathcal{I}_s[\varepsilon]:=a^{-s}\mathsf{P}_s\varepsilon=C_s .
}
\end{equation}
Thus the five constants \(\mathcal I_s\) are the invariant amplitudes of the round valuation-profile branches. Explicitly,
\begin{align}
\mathcal{I}_1
&=
\frac{a^{-1}}{24}\,
\Dil(\Dil+1)(\Dil+2)(\Dil+3)\varepsilon,
\\
\mathcal{I}_0
&=
-\frac{1}{6}\,
(\Dil-1)(\Dil+1)(\Dil+2)(\Dil+3)\varepsilon,
\\
\mathcal{I}_{-1}
&=
\frac{a}{4}\,
(\Dil-1)\Dil(\Dil+2)(\Dil+3)\varepsilon,
\\
\mathcal{I}_{-2}
&=
-\frac{a^2}{6}\,
(\Dil-1)\Dil(\Dil+1)(\Dil+3)\varepsilon,
\\
\mathcal{I}_{-3}
&=
\frac{a^3}{24}\,
(\Dil-1)\Dil(\Dil+1)(\Dil+2)\varepsilon.
\end{align}

\begin{corollary}[Boundary-only obstruction]
A density belongs to the boundary-only round valuation profile if and only if
\begin{equation}
\boxed{
\Dil(\Dil+1)(\Dil+2)(\Dil+3)\varepsilon=0 .
}
\end{equation}
Equivalently, inside the full valuation class the bulk obstruction is the invariant
\begin{equation}
\boxed{
\mathcal{I}_{\rm bulk}:=\mathcal{I}_1=\frac{\gamma A_0}{4}.
}
\end{equation}
The boundary-only restriction is precisely
\begin{equation}
\mathcal{I}_{\rm bulk}=0.
\end{equation}
\end{corollary}

\begin{corollary}[Minimal endpoint density obstruction]
A density belongs to the minimal endpoint form
\begin{equation}
\varepsilon(a)=\varepsilon_q+\varepsilon_w a^{-3}
\end{equation}
if and only if
\begin{equation}
\boxed{
\Dil(\Dil+3)\varepsilon=0 .
}
\end{equation}
Inside the full round valuation profile this is equivalent to the three obstructions
\begin{equation}
\mathcal{I}_1=0,
\qquad
\mathcal{I}_{-1}=0,
\qquad
\mathcal{I}_{-2}=0.
\end{equation}
The surviving invariant amplitudes are
\begin{equation}
\varepsilon_q=\mathcal{I}_0,
\qquad
\varepsilon_w=\mathcal{I}_{-3}.
\end{equation}
The identification of this density with an endpoint functional
\begin{equation}
E_{\rm end}=\sigma V_3+\mu Q
\end{equation}
requires the separate minimal endpoint assumptions.
\end{corollary}

\begin{remark}[Invariant content of the equation of state]
The equation-of-state values are not independent assumptions. They are the images of the dilation eigenvalues under
\begin{equation}
s\longmapsto w_s=-1-\frac{s}{3}.
\end{equation}
Hence
\begin{equation}
s=1,0,-1,-2,-3
\quad\Longleftrightarrow\quad
w=-\frac43,-1,-\frac23,-\frac13,0.
\end{equation}
The finite valuation spectrum is therefore a spectral statement for \(\Dil\), not merely a list of fitted powers.
\end{remark}

\begin{closurecondition}[Boundary-only valuation restriction]
Assume that the physical filling energy is required to be expressible entirely as a local valuation functional on the spatial representative \(\Sthree_A\), viewed as the boundary of the auxiliary filling \(\Bfour_A\),
\begin{equation}
\partial\Bfour_A=\Sthree_A,
\end{equation}
built from the boundary curvature measures of the homothetic convex filling. The boundary-only class used here means local scalar densities built only from the induced boundary measure and the elementary symmetric curvatures \(e_j(K)\), with coefficients independent of \(A\), and not boundary integrals involving the support function \(x\cdot n\) or other ambient-position data. Then the admissible valuation sector excludes the bulk intrinsic-volume term:
\begin{equation}
\gamma=0.
\end{equation}
It leaves the four boundary curvature slots
\begin{equation}
w=-1,\ -\frac{2}{3},\ -\frac{1}{3},\ 0.
\end{equation}
\end{closurecondition}
\begin{proof}
In the full valuation class the term
\begin{equation}
\gamma V_4(A)=\frac{\pi^2}{2}\gamma A^4
\end{equation}
is the four-dimensional bulk intrinsic-volume valuation. The boundary curvature valuations have homothetic powers
\begin{equation}
\int_{\Sthree_A}e_j(K)\,\dd\Sigma\propto A^{3-j},
\qquad j=0,1,2,3.
\end{equation}
Therefore the bulk \(A^4\) branch is not a member of the boundary curvature valuation sector. Under the boundary-only restriction, its coefficient must be set to zero. This proves \(\gamma=0\) relative to the stated boundary-only closure assumption, not from topology alone.
\end{proof}


The boundary-only restriction above is an independent geometric filter. Inside the explicitly stated minimal endpoint class, the following theorem gives a separate conditional endpoint classification. It is not a theorem from topology alone; it is conditionally unique only inside that class.

\begin{theorem}[Minimal endpoint classification]
Let
\begin{equation}
\Sigma=\partial\Bfour\simeq \Sthree .
\end{equation}
Assume that the endpoint sector is restricted to functionals satisfying the following conditions:
\begin{enumerate}
\item locality on the endpoint boundary \(\Sigma\);
\item invariance under \({\rm Diff}(\Sigma)\);
\item no local propagating endpoint degrees of freedom;
\item no duplication of gravitational curvature terms such as
\(\int_\Sigma \sqrt h\,R^{(3)}\), extrinsic-curvature actions, or higher-curvature endpoint terms;
\item no fixed external geometric structures, background connections, nonlocal spectral invariants, or additional geometric or dynamical length scales entering the functional dependence on \(A\) or on the endpoint fields; the overall dimensional coefficients remain free;
\item additivity under disjoint unions;
\item dependence on any internal field only through a topological charge.
\end{enumerate}
The additivity condition is understood on the extended class of finite disjoint unions of oriented endpoint components, and is then restricted back to the connected case \(\Sigma\simeq\Sthree\). Topological constants depending only on the number of connected components are excluded by the locality assumption and, after restriction back to connected \(\Sthree\), would in any case be an additional nonlocal topological input outside the minimal endpoint class.
Then the minimal endpoint energy has the form
\begin{equation}
E_{\rm end}(A,Q)
=\sigma\,\operatorname{Vol}(\Sthree_A)+\mu Q,
\end{equation}
where
\begin{equation}
Q[\Phi]=\int_{\Sthree}\Phi^*\omega_3\in\mathbb{Z}
\end{equation}
is the degree of a map
\begin{equation}
\Phi:\Sthree\longrightarrow \Sthree
\end{equation}
with \(\omega_3\) normalized by
\begin{equation}
\int_{\Sthree_{\rm unit}}\omega_3=1 .
\end{equation}
Thus
\begin{equation}
E_{\rm end}(A,Q)
=2\pi^2\sigma A^3+\mu Q .
\end{equation}
The induced endpoint density is
\begin{equation}
\epsD(A)
=\sigma+\frac{\mu Q}{2\pi^2A^3}.
\end{equation}
Equivalently, for \(A=A_0a\),
\begin{equation}
\epsD(a)
=\varepsilon_q+\varepsilon_w a^{-3},
\qquad
\varepsilon_q=\sigma,
\qquad
\varepsilon_w=\frac{\mu Q}{2\pi^2A_0^3}.
\end{equation}
The pressure obtained from
\begin{equation}
p_D=-\left(\frac{\partial E_{\rm end}}{\partial V_3}\right)_Q
\end{equation}
is
\begin{equation}
p_D=-\sigma=-\varepsilon_q .
\end{equation}
Hence the minimal endpoint sector is precisely the vacuum-dust background sector.
\end{theorem}
\begin{proof}
Within the declared local \({\rm Diff}(\Sigma)\)-invariant class and after excluding curvature, derivative, embedding-dependent, and background-structure terms, the only derivative-free metric scalar density is a constant multiple of the Riemannian volume density:
\begin{equation}
\int_{\Sigma}\dd V_3 .
\end{equation}
For the round representative this gives
\begin{equation}
\operatorname{Vol}(\Sthree_A)=2\pi^2A^3 .
\end{equation}
Its contribution is
\begin{equation}
E_q=\sigma V_3,
\end{equation}
so that
\begin{equation}
\varepsilon_q=\frac{E_q}{V_3}=\sigma,
\qquad
p_q=-\frac{\partial E_q}{\partial V_3}=-\sigma .
\end{equation}
By assumption, internal endpoint fields may enter only through a topological charge. Since
\begin{equation}
H^3(\Sthree;\mathbb{Z})\simeq\mathbb{Z},
\end{equation}
a map \(\Phi:\Sthree\to\Sthree\) gives the degree
\begin{equation}
Q[\Phi]=\int_{\Sthree}\Phi^*\omega_3\in\mathbb{Z}.
\end{equation}
This quantity is invariant under smooth deformations of \(\Phi\) inside a fixed homotopy class and is independent of the metric scale \(A\). Additivity under disjoint unions requires a linear dependence on the charge:
\begin{equation}
f(Q_1+Q_2)=f(Q_1)+f(Q_2),
\end{equation}
hence
\begin{equation}
f(Q)=\mu Q
\end{equation}
up to the choice of normalization of \(Q\). Therefore
\begin{equation}
E_w=\mu Q
\end{equation}
is independent of \(A\). Dividing by
\begin{equation}
V_3=2\pi^2A^3
\end{equation}
gives
\begin{equation}
\varepsilon_w(A)=\frac{\mu Q}{2\pi^2A^3}.
\end{equation}
At fixed topological charge,
\begin{equation}
p_w=-\left(\frac{\partial E_w}{\partial V_3}\right)_Q=0.
\end{equation}
All remaining local metric scalar terms contain intrinsic curvature, extrinsic curvature, higher derivatives, embedding data, fixed background structures, nonlocal spectral data, or additional geometric or dynamical length scales in the functional dependence. They either duplicate gravitational curvature terms, introduce additional endpoint dynamics, or leave the declared minimal endpoint class. The coefficients \(\sigma\) and \(\mu\) remain free dimensional amplitudes rather than additional functional dependence. Thus, inside the stated class, the endpoint functional is exhausted by
\begin{equation}
E_{\rm end}(A,Q)=\sigma V_3+\mu Q .
\end{equation}
\end{proof}

\subsection{Pressure Kernel and the Dust Branch}
The work-pressure map has a nontrivial kernel. This kernel is exactly the structural origin of the dust scaling in the minimal endpoint sector.

\begin{theorem}[Kernel of the work-pressure map]
Let \(E(A)\in C^1(J)\) on an interval \(J\subset(0,\infty)\), and let
\begin{equation}
\varepsilon(A)=\frac{E(A)}{V_3(A)},
\qquad
p(A)=-\frac{\dd E}{\dd V_3}.
\end{equation}
Then
\begin{equation}
p(A)=0
\end{equation}
on \(J\) if and only if
\begin{equation}
E(A)=E_0
\end{equation}
is constant on \(J\). In that case
\begin{equation}
\varepsilon(A)=\frac{E_0}{2\pi^2A^3}
=\varepsilon_w a^{-3}.
\end{equation}
Thus the homogeneous pressureless branch is precisely the image, under division by \(V_3\), of the constant-energy kernel of the work-pressure map.
\end{theorem}
\begin{proof}
Since
\begin{equation}
p=-\frac{\dd E}{\dd V_3}
=-\frac{\dd E/\dd A}{\dd V_3/\dd A},
\end{equation}
and since
\begin{equation}
\frac{\dd V_3}{\dd A}=6\pi^2A^2\neq0
\end{equation}
for \(A>0\), the condition \(p=0\) is equivalent to
\begin{equation}
\frac{\dd E}{\dd A}=0.
\end{equation}
Hence \(E(A)=E_0\). Division by \(V_3=2\pi^2A^3\) gives the dust scaling.
\end{proof}

\begin{corollary}[Endpoint normal form from volume plus pressure kernel]
The minimal endpoint energy
\begin{equation}
E_{\rm end}(A,Q)=\sigma V_3(A)+\mu Q
\end{equation}
is the direct sum of two distinguished work-map pieces:
\begin{equation}
\sigma V_3(A)
\quad\text{and}\quad
\mu Q.
\end{equation}
The first is the volume branch and gives
\begin{equation}
\varepsilon_q=\sigma,
\qquad
p_q=-\sigma.
\end{equation}
At fixed topological charge, the second lies in the pressure kernel and gives
\begin{equation}
\varepsilon_w(a)=\frac{\mu Q}{2\pi^2A_0^3}a^{-3},
\qquad
p_w=0.
\end{equation}
Therefore the endpoint vacuum-dust form is the normal form of the endpoint energy under the decomposition
\begin{equation}
E_{\rm end}
\in
\operatorname{span}\{V_3\}
\oplus
\ker\!\left(E\mapsto-\frac{\dd E}{\dd V_3}\right),
\end{equation}
with the second summand restricted by the endpoint classification to the additive topological charge term.
\end{corollary}

\begin{remark}[Scope of the endpoint classification]
This theorem is not a theorem from topology alone. It is a conditionally unique endpoint reduction inside the minimal endpoint class defined by locality, boundary diffeomorphism invariance, additivity, absence of local endpoint degrees of freedom, exclusion of gravitational curvature duplication, and absence of fixed external geometric structures, background connections, nonlocal spectral invariants, or additional geometric or dynamical length scales entering the functional dependence. The overall dimensional coefficients remain free. If curvature terms, gradient terms for \(\Phi\), embedding-dependent terms, nonlocal functionals, background connections, spectral invariants, or additional functional scales are admitted, the conditional uniqueness statement no longer holds.
\end{remark}

\begin{remark}[Relation to valuation coefficients]
The valuation endpoint notation
\begin{equation}
\alpha_0\int_{\Sthree_A}\dd\Sigma+
\alpha_3\int_{\Sthree_A}\det(K^a{}_b)\,\dd\Sigma
\end{equation}
contains the same two homothetic powers. In the normalized topological-charge convention used in the minimal endpoint classification,
\begin{equation}
2\pi^2\alpha_0A^3+2\pi^2\alpha_3
=2\pi^2\sigma A^3+\mu Q,
\end{equation}
so that
\begin{equation}
\alpha_0=\sigma,
\qquad
\alpha_3=\frac{\mu Q}{2\pi^2}.
\end{equation}
For the Gauss-map endpoint of the convex round filling, the normalized degree is \(Q=1\). The identification of the valuation constant slot with the endpoint charge is a matching of homothetic powers and, in the convex round case, of the Gauss-map degree. It is not an admission of an independent extrinsic-curvature action into the minimal endpoint class. More generally, an internal map \(\Phi:\Sthree\to\Sthree\) supplies the integer charge. The topological charge selects the dust background scaling; it is not by itself a cold-dark-matter perturbation theory.
\end{remark}

\begin{remark}[Why the endpoint terms are distinguished]
The derivative-free metric scalar density is the constant zeroth-order density on the boundary, which is integrated into \(\sigma V_3\). The topological endpoint contribution is controlled by the integer class \(H^3(\Sthree;\mathbb{Z})\simeq\mathbb{Z}\), represented by \(Q[\Phi]\). In this sense
\begin{equation}
H^0(\Sthree)\leadsto \sigma V_3,
\qquad
H^3(\Sthree;\mathbb{Z})\leadsto \mu Q.
\end{equation}
This explains why the minimal endpoint functional contains precisely the \(A^3\) and \(A^0\) powers before any covariant realization of dust perturbations is imposed.
\end{remark}

\begin{remark}[Gauss--Kronecker term and Gauss-map degree]
The endpoint invariant
\begin{equation}
\int_{\Sthree_A}\det(K^a{}_b)\,\dd\Sigma
\end{equation}
is not an Euler characteristic of \(\Sthree\). Since \(\Sthree\) is odd-dimensional,
\begin{equation}
\chi(\Sthree)=0.
\end{equation}
For an oriented closed hypersurface in the auxiliary Euclidean space \(\mathbb{R}^4\), the Gauss--Kronecker integral is instead the Jacobian degree of the Gauss map
\begin{equation}
\nu:\Sthree_A\longrightarrow \Sthree_{\rm unit}.
\end{equation}
With the chosen outward orientation,
\begin{equation}
\int_{\Sthree_A}\det(K^a{}_b)\,\dd\Sigma
=\deg(\nu)\,\vol(\Sthree_{\rm unit}).
\end{equation}
For the convex round filling, \(\deg(\nu)=1\), and therefore
\begin{equation}
\int_{\Sthree_A}\det(K^a{}_b)\,\dd\Sigma=2\pi^2.
\end{equation}
Thus the top endpoint term is a Gauss-map degree invariant in the round convex case, not a Gauss--Bonnet--Chern Euler term. In the intrinsic-volume language, this top boundary curvature integral corresponds, up to the normalization absorbed into \(\alpha_3\), to the zeroth intrinsic volume of the convex body. The sign of the degree depends on the orientation convention for the normal; throughout this paper the outward normal of the convex filling is used.
\end{remark}

\begin{nogotheorem}[No amplitude from topology alone]
The topology of the carrier \(\Sthree\), the ratio \(V_4/V_3=A/4\), the full valuation spectrum, the boundary-only restriction, and the minimal endpoint classification do not determine the numerical amplitudes \(\sigma\) and \(\mu\).
\end{nogotheorem}

\begin{proof}
The two minimal endpoint terms are
\begin{equation}
V_3(A)=\int_{\Sthree_A}\dd V_3,
\qquad
Q[\Phi]=\int_{\Sthree}\Phi^*\omega_3 .
\end{equation}
Multiplying either term by an arbitrary constant with the required physical dimension preserves the endpoint assumptions: locality on the boundary, boundary reparameterization invariance, additivity, absence of local endpoint degrees of freedom, absence of fixed external structures, and dependence on internal fields only through the topological charge. Thus the minimal endpoint class fixes the powers \(A^3\) and \(A^0\), but not the dimensional coefficients multiplying them. The Gauss--Kronecker endpoint is one geometric realization of the constant \(A^0\) valuation power in the round convex filling, while \(Q[\Phi]\) is the normalized topological-charge realization used in the minimal endpoint class. A numerical amplitude requires an additional input: a quantum scale, a matter normalization, a compensation law, an observational calibration, or another independent dimensional principle.
\end{proof}

\begin{corollary}[No amplitude from endpoint linearity]
The same no-go applies to every linear combination of the endpoint invariants. The minimal endpoint classification fixes the two powers \(A^3\) and \(A^0\), but the coefficients multiplying them remain free dimensional constants until an independent physical or observational input is supplied.
\end{corollary}

\begin{corollary}[Exact scope of predictivity]
The full valuation theorem gives the five-slot homothetic spectrum
\begin{equation}
\epsD(A)=\frac{\gamma A}{4}+\alpha_0+\frac{3\alpha_1}{A}+\frac{3\alpha_2}{A^2}+\frac{\alpha_3}{A^3}.
\end{equation}
The boundary-only restriction \(\gamma=0\) removes the bulk branch relative to the explicitly stated boundary-only valuation restriction. Inside the stated minimal endpoint class one obtains the functional form
\begin{equation}
\epsD(a)=\varepsilon_q+\varepsilon_w a^{-3},
\qquad
\pD=-\varepsilon_q,
\end{equation}
but not the numerical values of \(\varepsilon_q\) and \(\varepsilon_w\). Thus the minimal endpoint class reduces the arbitrary background function \(\epsD(A)\) to two constants after the stated assumptions. It does not determine those constants by topology alone.
\end{corollary}

\section{Minimal Covariant Effective Completion of the Selected Sector}
The minimal endpoint classification identifies the vacuum-dust background form inside the stated endpoint class. The boundary-only valuation restriction supplies an independent geometric route to elimination of the bulk branch. A covariant four-dimensional completion is still required in order to define the selected sector explicitly. The minimal completion matching the selected sector is the standard vacuum plus conserved dust current, in the standard variational-fluid sense \cite{Brown1993Fluid}.

Let
\begin{equation}
S_{\rm eff}=S_{\rm grav}+S_{\rm vis}+S_q+S_w,
\end{equation}
with
\begin{equation}
S_q=-\int_{\M}\sqrt{-g}\,\varepsilon_q\,\dd^4x .
\end{equation}
Then
\begin{equation}
T^q_{\mu\nu}=-\varepsilon_q g_{\mu\nu}.
\end{equation}

For the dust branch introduce three comoving scalar labels \(Z^A\), \(A=1,2,3\), and the identically conserved particle-number current density. Here \(\epsilon^{\mu\nu\rho\sigma}\) denotes the alternating symbol, so \(J^\mu\) is a contravariant vector density of weight \(+1\):
\begin{equation}
J^\mu=\frac{1}{3!}\epsilon^{\mu\nu\rho\sigma}\epsilon_{ABC}
\partial_\nu Z^A\partial_\rho Z^B\partial_\sigma Z^C,
\qquad
\partial_\mu J^\mu=0.
\end{equation}
Define
\begin{equation}
n=\frac{\sqrt{-g_{\mu\nu}J^\mu J^\nu}}{\sqrt{-g}},
\qquad
u^\mu=\frac{J^\mu}{\sqrt{-g}\,n},
\qquad
u^\mu u_\mu=-1.
\end{equation}
The ordinary conservation law \(\partial_\mu J^\mu=0\) is equivalent to
\begin{equation}
\nabla_\mu(n u^\mu)=0.
\end{equation}
The dust action is
\begin{equation}
S_w=-m\int_{\M}\sqrt{-g}\,n\,\dd^4x .
\end{equation}
Its stress tensor is
\begin{equation}
T^w_{\mu\nu}=\rho_w u_\mu u_\nu,
\qquad
\rho_w=mn.
\end{equation}
Together with \(\nabla_\mu(nu^\mu)=0\), conservation of \(T^w_{\mu\nu}=\rho_w u_\mu u_\nu\) splits into continuity and geodesic motion. In the dust branch this gives
\begin{equation}
u^\mu\nabla_\mu u^\nu=0.
\end{equation}

On the closed FLRW background, \(n=n_0a^{-3}\). Hence the constant dust coefficient in \(\epsD=\varepsilon_q+\varepsilon_w a^{-3}\) is
\begin{equation}
\varepsilon_w=mn_0.
\end{equation}
Matching the homogeneous dust energy to the endpoint constant gives
\begin{equation}
E_w=V_3(A)\,mn_0a^{-3}=2\pi^2A_0^3mn_0=\mu Q,
\end{equation}
and therefore
\begin{equation}
\varepsilon_w=mn_0=\frac{\mu Q}{2\pi^2A_0^3}.
\end{equation}
If the normalized topological charge is identified with the comoving particle number,
\begin{equation}
Q=2\pi^2A_0^3n_0,
\end{equation}
then \(\mu=m\) in this normalization.
The total selected residual tensor is therefore
\begin{equation}
(T^D)^\mu{}_{\nu}
=\diag\bigl(-\varepsilon_q-\varepsilon_w a^{-3},
-\varepsilon_q,-\varepsilon_q,-\varepsilon_q\bigr),
\end{equation}
so that
\begin{equation}
\epsD=\varepsilon_q+\varepsilon_w a^{-3},
\qquad
\pD=-\varepsilon_q.
\end{equation}

\begin{theorem}[Cold pressureless behavior in the minimal conserved-current realization]
For the minimal conserved-current completion of the \(A^{-3}\) branch,
\begin{equation}
c_s^2=0,
\qquad
\pi_{\mu\nu}=0,
\qquad
u^\mu\nabla_\mu u^\nu=0.
\end{equation}
In the minimal conserved-current realization, the endpoint \(A^{-3}\) branch has the stress tensor and perturbative closure of cold pressureless matter.
\end{theorem}

\begin{remark}[Effective matching]
This completion is a four-dimensional effective matching completion. It does not prove that the microscopic origin of the dust current is the Gauss-Kronecker boundary invariant. Rather, the minimal endpoint classification contains the homogeneous \(A^{-3}\) branch, and the conserved-current action provides a covariant effective realization with the required perturbative properties.
\end{remark}

\begin{remark}[Existence, not uniqueness]
This construction proves the existence of a covariant effective completion matching the endpoint-selected background sector. It does not prove uniqueness. Other completions with additional fields, nonminimal couplings, or boundary degrees of freedom would have to reproduce the same background tensor and pass the perturbative criteria independently.
\end{remark}

\section{Topological Charge and Dimensional Amplitudes}
The minimal endpoint classification contains a normalized topological charge \(Q\), but it still does not create a dimensional mass scale. A stronger future completion may identify this endpoint charge with the winding number of a genuine topological field.

Let
\begin{equation}
U:\Sigma\simeq\Sthree\longrightarrow SU(2)\simeq\Sthree
\end{equation}
be a smooth map, and let
\begin{equation}
L=U^{-1}\dd U.
\end{equation}
With the chosen trace normalization and orientation, the winding number is
\begin{equation}
N[U]=
\frac{1}{24\pi^2}
\int_{\Sthree}
\operatorname{tr}(L\wedge L\wedge L)
\in\mathbb{Z},
\end{equation}
reflecting \(\pi_3(SU(2))\simeq\mathbb{Z}\). The trace normalization is chosen so that the identity map \(\Sthree\simeq SU(2)\to SU(2)\) has \(N[U]=1\). This is the standard topological-charge normalization used in Skyrme-type constructions and current-algebra topological actions \cite{Skyrme1961,Witten1983CurrentAlgebra}.

The classification theorem uses the oriented charge \(Q\in\mathbb{Z}\). For a physical dust interpretation one may subsequently restrict to a fixed-orientation superselection sector \(Q\ge0\), in which the positive particle number is \(N=Q\). If both orientations are allowed, replacing \(Q\) by \(|Q|\) is a further physical restriction and is not the same additive valuation. In the fixed-orientation sector one may write
\begin{equation}
Q=N=N[U]\ge0,
\qquad N\in\mathbb{Z}_{\ge0}.
\end{equation}
If the dust branch is matched by
\begin{equation}
E_w=mN,
\end{equation}
then the endpoint term \(E_w=\mu Q\) gives \(\mu=m\) in this normalization. Consistency with the conserved-current dust normalization gives
\begin{equation}
Q=N=2\pi^2A_0^3 n_0,
\qquad
\varepsilon_w=mn_0=\frac{mN}{2\pi^2A_0^3}
=\frac{\mu Q}{2\pi^2A_0^3}.
\end{equation}
Thus this optional topological completion would impose a quantization condition on the comoving particle number \(2\pi^2A_0^3n_0\). Without that additional identification, \(n_0\) remains the continuous normalization of the effective dust current.

Topology may quantize the dimensionless number \(Q\) or \(N\), but it does not generate the mass-per-charge scale \(\mu\). The dimensional dust amplitude is controlled by the product \(\mu Q\): topology supplies the integer factor, while physics must supply the scale.

\begin{remark}[Maurer--Cartan forms versus dynamical fields]
The Maurer--Cartan forms of the spatial \(\Sthree\) are geometric coframes, not automatically independent four-dimensional matter fields. A covariant topological completion must introduce a field
\begin{equation}
U(x):\M\longrightarrow SU(2),
\qquad
L_\mu=U^{-1}\partial_\mu U,
\end{equation}
and may then define the topological current density
\begin{equation}
J^\mu_{\rm top}
=
\frac{1}{24\pi^2}
\epsilon^{\mu\nu\rho\sigma}
\operatorname{tr}(L_\nu L_\rho L_\sigma),
\qquad
\partial_\mu J^\mu_{\rm top}=0.
\end{equation}
This current supplies a conserved charge. It is not, by itself, a dust stress tensor. To reproduce the minimal endpoint dust sector, the completed action must still yield
\begin{equation}
T_{\mu\nu}=\rho u_\mu u_\nu,
\qquad
c_s^2=0,
\qquad
\pi_{\mu\nu}=0,
\end{equation}
in the relevant effective limit.
\end{remark}


\begin{theorem}[Conditional pressureless action for a topological current]
Let a future completion introduce a smooth map
\begin{equation}
U:\M\longrightarrow SU(2),
\qquad
L_\mu=U^{-1}\partial_\mu U.
\end{equation}
Let \(\vartheta=g^{-1}\dd g\) denote the Maurer--Cartan form on \(SU(2)\), and let
\begin{equation}
\omega_{SU(2)}
=
\frac{1}{24\pi^2}
\operatorname{tr}(\vartheta\wedge\vartheta\wedge\vartheta)
\end{equation}
be the normalized closed volume three-form on \(SU(2)\). Then
\begin{equation}
U^*\omega_{SU(2)}
=
\frac{1}{24\pi^2}
\operatorname{tr}(L\wedge L\wedge L),
\qquad
L=U^{-1}\dd U.
\end{equation}
Since \(\dd\omega_{SU(2)}=0\), naturality gives
\begin{equation}
\dd\,U^*\omega_{SU(2)}=U^*(\dd\omega_{SU(2)})=0.
\end{equation}
The dual topological current density is
\begin{equation}
J^\mu_{\rm top}
=
\frac{1}{24\pi^2}
\epsilon^{\mu\nu\rho\sigma}
\operatorname{tr}(L_\nu L_\rho L_\sigma),
\end{equation}
and satisfies
\begin{equation}
\partial_\mu J^\mu_{\rm top}=0.
\end{equation}
If, in addition, the effective dust action is chosen to be
\begin{equation}
S_{\rm top\text{-}dust}
=-m\int_{\M}\sqrt{-g}\,
\frac{\sqrt{-g_{\mu\nu}J^\mu_{\rm top}J^\nu_{\rm top}}}{\sqrt{-g}}\,
\dd^4x,
\end{equation}
on the sector where \(J^\mu_{\rm top}\) is timelike and future-directed, then the resulting stress tensor is
\begin{equation}
T_{\mu\nu}=\rho u_\mu u_\nu,
\qquad
\rho=mn,
\qquad
u^\mu=\frac{J^\mu_{\rm top}}{\sqrt{-g}\,n},
\end{equation}
where
\begin{equation}
n=\frac{\sqrt{-g_{\mu\nu}J^\mu_{\rm top}J^\nu_{\rm top}}}{\sqrt{-g}}.
\end{equation}
Consequently this chosen effective completion has
\begin{equation}
c_s^2=0,
\qquad
\pi_{\mu\nu}=0.
\end{equation}
\end{theorem}
\begin{proof}
The conservation identity follows from the closure of the target-space volume three-form and naturality of pullback, \(\dd\,U^*\omega_{SU(2)}=U^*(\dd\omega_{SU(2)})=0\). With the stated action, \(J^\mu_{\rm top}\) plays the same role as the conserved particle-number current density in the standard pressureless-current action. Metric variation at fixed current density gives
\begin{equation}
T_{\mu\nu}=mn\,u_\mu u_\nu.
\end{equation}
On solutions of the \(U\)-field equations, diffeomorphism invariance gives
\begin{equation}
\nabla_\mu T^{\mu\nu}=0.
\end{equation}
Together with
\begin{equation}
\nabla_\mu(nu^\mu)=0,
\end{equation}
this implies
\begin{equation}
u^\mu\nabla_\mu u^\nu=0.
\end{equation}
A pressureless current action has no rest-frame pressure perturbation and no anisotropic stress in this effective sector, which gives the stated \(c_s^2\) and \(\pi_{\mu\nu}\). The result is conditional on the action matching; the topological current alone supplies only a conserved charge.
\end{proof}

\begin{remark}[Scope of the topological-current realization]
The action gives a dust stress tensor on the sector where the topological current is timelike, but it restricts the allowed dust configurations to those representable by a map \(U:\M\to SU(2)\). Thus it is a special topological realization of the conserved-current dust action, not the full Brown dust theory.
\end{remark}

\begin{remark}[Vacuum amplitude]
The same topological charge does not determine the vacuum coefficient \(\sigma\). In units \(\hbar=c=1\), a vacuum scale would have the form
\begin{equation}
\sigma=C_\sigma\Lambda_*^4,
\end{equation}
or it must be fixed observationally, for example by \(\sigma=\rho_{\Lambda,0}\). The dimensionless constant \(C_\sigma\) and the scale \(\Lambda_*\) are not fixed by \(\Sthree=\partial\Bfour\), the valuation basis, or the minimal endpoint classification alone.
\end{remark}

\section{Minisuperspace Contribution to the Total Action}
The full background action must have the form
\begin{equation}
S_{\rm tot}=S_{\rm grav}[A,N]+S_{\rm vis}[A,N,\psi]+S_{\rm fill}[A,N,\xi].
\end{equation}
The filling contribution is
\begin{equation}
S_{\rm fill}=-\int \dd t\,N(t)\Efour(A(t),\xi(t)).
\end{equation}
These reduced variational formulae hold in the sector where \(\Efour=\Efour(A,\xi)\) has no independent dependence on \(N\), \(\dot A\), \(\dot\xi\), curvature scalars, or extrinsic-curvature boundary data. Such dependences belong to the covariant completion problem.
Its variations are
\begin{equation}
\frac{\delta S_{\rm fill}}{\delta N(t)}=-\Efour,
\qquad
\frac{\delta S_{\rm fill}}{\delta A(t)}=-N\left.\frac{\partial\Efour}{\partial A}\right|_{\xi}.
\end{equation}
Hence
\begin{equation}
\epsD=-\frac{1}{V_3}\frac{\delta S_{\rm fill}}{\delta N},
\qquad
\pD=\frac{A}{3NV_3}\frac{\delta S_{\rm fill}}{\delta A}.
\end{equation}

\begin{remark}
The functional \(S_{\rm fill}\) determines the contribution of the filling sector to the full equations. It cannot be varied as a self-contained universe: equations of motion follow from \(S_{\rm tot}\), not from imposing \(\delta S_{\rm fill}=0\) separately. What has been constructed is a consistent minisuperspace contribution, not a complete local four-dimensional action. The distinction between a reduced action and a completed gravitational variational principle is standard in canonical and boundary-action treatments \cite{DeWitt1967,GibbonsHawking1977}.
\end{remark}

\section{Reconstructed, Constructed, and Physical Tensors}
For a specified metric and visible sector, define the reconstructed target residual
\begin{equation}
T^{\rm rec}_{\mu\nu}:=\frac{c^4}{8\pi G}G_{\mu\nu}-T^{\rm vis}_{\mu\nu}.
\end{equation}
This is a covariant reconstruction, not a microscopic derivation.

Let \(T^{\rm fill}_{\mu\nu}\) denote the tensor independently obtained from a covariant completion of the filling functional. On the FLRW background, its mixed components must reduce to
\begin{equation}
(T^{\rm fill})^{\mu}{}_{\nu}=\diag(-\epsD,\pD,\pD,\pD).
\end{equation}
The substantive matching condition is
\begin{equation}
T^{\rm fill}_{\mu\nu}=T^{\rm rec}_{\mu\nu}.
\end{equation}
Only after this has been satisfied are both objects denoted by the physical dark tensor \(T^D_{\mu\nu}\).

The Bianchi identities imply
\begin{equation}
\nabla_\mu(T^{\rm rec})^{\mu\nu}=-\nabla_\mu(T^{\rm vis})^{\mu\nu}.
\end{equation}
If \(\nabla_\mu(T^{\rm vis})^{\mu\nu}=J^\nu\), then \(\nabla_\mu(T^{\rm rec})^{\mu\nu}=-J^\nu\). For \(J^\nu=0\), the reconstructed residual is conserved. This condition does not replace the equations for the internal variables of the filling sector.

\section{\texorpdfstring{$\Sfour$}{S4} Compensation by Gluing}
Consider the gluing
\begin{equation}
\Sfour=\Bfour_+\cup_f\Bfour_-,
\end{equation}
where \(f:\Sthree_+\to\Sthree_-\) identifies the common boundaries. The following statement is only a boundary-orientation lemma; it does not prove cancellation of energies, actions, stresses, amplitudes, or equations of motion.

\begin{proposition}[Boundary-form cancellation, not dynamical compensation]
Let \(\omega_\pm\) be boundary forms satisfying
\begin{equation}
\iota_+^*\omega_+=f^*\iota_-^*\omega_-.
\end{equation}
Then, with opposite induced orientations,
\begin{equation}
\int_{\partial\Bfour_+}\omega_+ + \int_{\partial\Bfour_-}\omega_-=0.
\end{equation}
\end{proposition}
\begin{proof}
The gluing map \(f\) identifies the values of the forms on the common boundary, while the induced orientations supply opposite signs. The two contributions are therefore equal in magnitude and opposite in sign.
\end{proof}

\begin{remark}[Chern--Simons boundary actions and amplitude limitation]
For a boundary connection \(\mathcal{A}\), one may form the Chern--Simons three-form \cite{ChernSimons1974}
\begin{equation}
\operatorname{CS}_3(\mathcal{A})=
\operatorname{tr}\left(\mathcal{A}\wedge\dd\mathcal{A}
+\frac{2}{3}\mathcal{A}\wedge\mathcal{A}\wedge\mathcal{A}\right),
\end{equation}
with
\begin{equation}
\dd\operatorname{CS}_3(\mathcal{A})=\operatorname{tr}(F_{\mathcal{A}}\wedge F_{\mathcal{A}}).
\end{equation}
For a gluing \(\Sfour=\Bfour_+\cup_f\Bfour_-\), compatible boundary data and opposite induced orientations imply cancellation of the oriented boundary contributions up to the usual quantized large-gauge shift fixed by the normalization of the Chern--Simons coupling. This is a statement about oriented topological consistency of boundary data. It does not determine the dimensional amplitudes \(\sigma\) or \(\mu\), which require independent couplings or physical scales.
\end{remark}


\begin{nogotheorem}[No vacuum amplitude from Chern--Simons gluing alone]
A Chern--Simons boundary functional can supply an oriented topological phase or a quantized large-gauge shift, but it does not by itself determine the dimensionful vacuum energy density \(\sigma\).
\end{nogotheorem}
\begin{proof}
The Chern--Simons functional changes under a large gauge transformation by a quantized amount fixed by the normalization of the Chern--Simons coupling. This quantized shift is a statement about the boundary action or phase. The endpoint vacuum contribution is instead
\begin{equation}
E_q=\sigma V_3=2\pi^2\sigma A^3,
\end{equation}
where
\begin{equation}
[\sigma]=EL^{-3}.
\end{equation}
Equating \(E_q\) to a Chern--Simons shift requires an additional dimensionful coupling, time scale, cutoff, or Hamiltonian normalization. Therefore Chern--Simons gluing may constrain a topological sector or a phase, but it cannot determine \(\sigma\) without an independent dimensional input.
\end{proof}

Orientation supplies the sign but not equality of the boundary data. Functionals containing normal derivatives, extrinsic curvature, independent fields, or corner terms require additional matching conditions. The construction of a concrete \(\Sfour\) functional that selects the amplitudes \(\varepsilon_q\) and \(\varepsilon_w\) is therefore a separate completion condition.

\section{Suspension, \texorpdfstring{$SU(2)$}{SU(2)}, and Hopf Structure}
We retain the topological chain
\begin{equation}
S^0\longrightarrow S^1\longrightarrow S^2\longrightarrow \Sthree\longrightarrow \Sfour,
\end{equation}
with \(\Sfour\simeq \Sigma\Sthree\). The exact structural facts are
\begin{equation}
\Sthree\simeq SU(2),
\qquad
T\Sthree\simeq \Sthree\times\mathbb{R}^3,
\qquad
S^1\hookrightarrow \Sthree\to S^2.
\end{equation}
For the chosen normalization of the left-invariant Maurer-Cartan forms,
\begin{equation}
\dd\sigma^i=-\frac{1}{2}\epsilon^i{}_{jk}\sigma^j\wedge\sigma^k.
\end{equation}
Parallelizability and the Hopf fibration provide global geometric structures capable of carrying directions, phases, and flows \cite{Hopf1931}, but they do not by themselves determine an energy or the amplitude of the residual sector.

\begin{hypothesis}
Global Hopf modes may enter the internal variables \(\xi^I\) of the filling energy. Their physical imprint would then have to appear through \(\QD\), anisotropic stress, a surface charge, or the spectrum of global modes, rather than through a topological identification alone.
\end{hypothesis}

\section{Symmetries and Noether Identities}
Suppose a four-dimensional diffeomorphism-invariant completion \(S^{\rm eff}_{\rm fill}[g,\xi]\) exists and the \(\xi^I\) are scalar fields. The distinctions used here follow the standard variational and covariant phase-space separation of identities, currents, constraints, and charges \cite{Noether1918,LeeWald1990,IyerWald1994,BarnichBrandt2002}. We use the stress-tensor convention
\begin{equation}
T_{\mu\nu}:=-\frac{2}{\sqrt{-g}}\frac{\delta S}{\delta g^{\mu\nu}}.
\end{equation}
With conventions
\begin{equation}
\delta_\zeta \xi^I=\zeta^\mu\nabla_\mu\xi^I,
\qquad
\mathcal{E}_I:=\frac{1}{\sqrt{-g}}\frac{\delta S^{\rm eff}_{\rm fill}}{\delta\xi^I},
\end{equation}
diffeomorphism invariance gives
\begin{equation}
\nabla_\mu(T^{\rm fill})^{\mu}{}_{\nu}=\sum_I\mathcal{E}_I\nabla_\nu\xi^I.
\end{equation}
On shell, \(\mathcal{E}_I=0\), and therefore \(\nabla_\mu(T^{\rm fill})^\mu{}_{\nu}=0\). For tensor or gauge fields, the right-hand side is replaced by the contraction of the Euler-Lagrange expressions with the corresponding diffeomorphism generators. This is a variational identity, not an independent global current.

Four objects must be kept distinct: global current, gauge identity, constraint, surface charge. A superpotential belongs to the decomposition of a Noether current or a surface charge; it is not appended arbitrarily to the divergence of an energy-momentum tensor. A possible covariant boundary completion must likewise distinguish bulk currents from quasilocal surface charges.

\section{Light and a Causal Measure}
The radial null condition for the metric gives
\begin{equation}
\dd\chi=\frac{cN(t)}{A(t)}\dd t.
\end{equation}
Define the causal angular radius
\begin{equation}
\chi_c(t_i,t):=\min\left\{\pi,\int_{t_i}^t\frac{cN(t')}{A(t')}\,\dd t'\right\}.
\end{equation}
The causally accessible fraction of the three-volume is
\begin{equation}
I_c(t_i,t):=\frac{V_3(\chi_c,A)}{V_3(\pi,A)}
=\frac{\chi_c-\sin\chi_c\cos\chi_c}{\pi},
\qquad
0\le I_c\le1.
\end{equation}
The antipode is reached at accumulated conformal length \(\pi\), while a complete geodesic circuit requires \(2\pi\). Causal accessibility is not identified with entropic or quantum coherence.

\begin{hypothesis}
If \(\Efour=\Efour(A,I_c,\xi)\), then the causal variable contributes a definite exchange term,
\begin{equation}
Q_D^{(I_c)}=\frac{1}{V_3}\frac{\partial\Efour}{\partial I_c}\frac{\dd I_c}{\dd\tau}.
\end{equation}
\end{hypothesis}

\section{Radius Protection}
\begin{nogotheorem}
The topological class \(\Sthree\) alone provides neither a lower nor an upper bound on the radius \(A\).
\end{nogotheorem}
\begin{proof}
For every \(A>0\), the metric \(g_A=A^2\dd\Omega_3^2\) is a smooth round metric on the same manifold \(\Sthree\). Changing \(A\) does not change the topological class.
\end{proof}

Here \(A\) belongs to the metric representative \(g_A\), whereas the topological carrier remains \(\Sthree\) for every \(A>0\). The limits \(A\to0\) and \(A\to\infty\) are respectively a degenerate boundary and an asymptotic boundary of the homothetic metric family. They are not established physical solutions of the present construction, because the background map does not supply a completed dynamics selecting a trajectory \(A(\tau)\).

Topology alone therefore neither forbids nor predicts dynamical approach to either limit. Radius protection is an admissibility and dynamical problem. Possible additional inputs include a spectral or Planck-scale cutoff, an energetic threshold, a quasilocal balance law, finite causal completion, an infrared spectral condition, or an \(\Sfour\)-compensation law; none is derived from \(\Sthree\) topology alone.

\begin{remark}[Limits of controlled reconstruction]
At sufficiently small \(A\), Planck-scale curvature or spectral effects may mark a boundary of controlled reconstruction. Loss of reconstruction control is not equivalent to ontological destruction of the carrier. Likewise, a formal future extrapolation beyond the controlled domain is not a prediction of a Big Rip. Such a prediction would require a completed dynamics selecting a divergent evolution, together with its proper-time and physical-curvature behaviour.
\end{remark}

\section{Perturbations and Required Closure}
The background map does not determine perturbations automatically. The perturbative variables and closure requirements follow the standard cosmological perturbation framework \cite{KodamaSasaki1984,HalliwellHawking1985,MukhanovFeldmanBrandenberger1992,MalikWands2009}. A minimal set includes
\begin{equation}
\delta\rhoFour,
\qquad
\delta\epsD,
\qquad
\delta\pD,
\qquad
\delta\QD,
\end{equation}
metric scalar modes on \(\Sthree\), vector and tensor harmonics, and
\begin{equation}
c_s^2,
\qquad
\Gamma_D,
\qquad
\pi_{\mu\nu}.
\end{equation}
Here \(c_s^2\) is the effective sound speed in the chosen physical frame, \(\Gamma_D\) is the nonadiabatic perturbation, and \(\pi_{\mu\nu}\) is the anisotropic stress. Their determination requires a covariant completion or independent equations for the internal variables. On compact \(\Sthree\), a complete treatment must additionally respect the discrete scalar, vector, and tensor harmonic spectra.

For a generic filling functional, realizing a dust slot as cold dark matter requires at least
\begin{equation}
c_s^2=0,
\qquad
\pi_{\mu\nu}=0,
\end{equation}
together with geodesic motion of the effective flow. The conserved-current completion of the endpoint-selected \(A^{-3}\) branch satisfies these conditions. Non-minimal completions must prove them independently. The vacuum slot likewise requires a covariant condition eliminating an independent clustering perturbation; the background equality \(w=-1\) alone does not prove this.

\begin{closurecondition}[Perturbative completion beyond the minimal sector]
For extensions beyond the vacuum plus conserved-dust completion, construct the perturbed filling functional and derive the scalar, vector, and tensor equations on \(\Sthree\), including \(\delta\pD\), \(c_s^2\), \(\Gamma_D\), \(\pi_{\mu\nu}\), and the exchange law.
\end{closurecondition}



\section{Two Independent Mathematical Bases and Endpoint Classification}
The construction rests on two independent mathematical structures that restrict admissible residual-sector solutions at different levels: the valuation closure of the auxiliary filling and the compact harmonic spectrum of the spatial representative. The former constrains homogeneous background scaling laws, while the latter constrains the perturbative mode content. The minimal endpoint classification then selects the two endpoint powers inside a stated local boundary class. None of these structures fixes dimensional amplitudes or establishes observational superiority by itself.

\subsection{Valuation closure}
Under the Euclidean convex valuation hypothesis stated above, Hadwiger's classification restricts the auxiliary filling energy to the intrinsic-volume span \cite{Hadwiger1957,Schneider2014}. For the smooth round homothetic family this is represented by
\begin{equation}
\Efour^{\rm val}(A)
=\gamma V_4(A)
+\sum_{j=0}^{3}\alpha_j
\int_{\Sthree_A}e_j(K)\,\dd\Sigma .
\end{equation}
With the outward-normal convention,
\begin{equation}
K^a{}_b=A^{-1}\delta^a_b,
\qquad
e_j(K)=\binom{3}{j}A^{-j},
\end{equation}
and therefore
\begin{equation}
\epsD(A)
=\frac{\gamma A}{4}
+\alpha_0
+\frac{3\alpha_1}{A}
+\frac{3\alpha_2}{A^2}
+\frac{\alpha_3}{A^3}.
\end{equation}
The associated pure branches have
\begin{equation}
w=-\frac{4}{3},\ -1,\ -\frac{2}{3},\ -\frac{1}{3},\ 0.
\end{equation}
Thus arbitrary homogeneous functions \(\epsD(A)\) are excluded inside the stated valuation class. This is a finite-spectrum theorem relative to the Euclidean convex valuation hypothesis, not a prohibition on other dark-sector constructions outside that class.

\subsection{Compact harmonic spectrum}
The second structure is harmonic analysis on the compact round spatial representative. For curvature radius \(R_c\), the scalar Laplacian has the discrete spectrum \cite{HalliwellHawking1985,KodamaSasaki1984}
\begin{equation}
\operatorname{Spec}(-\nabla^2_{\Sthree})
=
\left\{
\frac{n(n+2)}{R_c^2}
\right\}_{n=0}^{\infty},
\qquad
d_n=(n+1)^2,
\end{equation}
where \(d_n\) is the degeneracy of the scalar eigenspace. The mode \(n=0\) is homogeneous. Non-constant scalar harmonics begin at \(n=1\), subject to the usual gauge and constraint qualifications of cosmological perturbation theory. The harmonic spectrum is geometric; the physical perturbation spectrum is obtained only after constraints and gauge removal. Consequently, compact \(\Sthree\) does not possess a continuum of arbitrarily small nonzero scalar wavenumbers, unlike non-compact \(\mathbb{R}^3\).

This discreteness supplies a possible infrared spectral cutoff when the curvature scale, primordial spectrum, and transfer functions make it observationally accessible. It does not by itself imply suppression of a particular CMB multipole. Any observable infrared signature depends on \(R_c\), the allowed scalar, vector, and tensor harmonics, transfer functions, initial conditions, and any quotient topology \(\Sthree/\Gamma\).

\subsection{Synthesis}
The two structures constrain different layers of the model:
\begin{equation}
\boxed{
\text{valuation closure on }\Bfour_A
\quad\Longrightarrow\quad
\text{finite background equation-of-state spectrum}
}
\end{equation}
and
\begin{equation}
\boxed{
\begin{gathered}
\text{compact harmonics on }\Sthree_A\\
\Longrightarrow\ \text{discrete perturbation spectrum and possible IR signatures}
\end{gathered}
}.
\end{equation}
The minimal endpoint classification adds a separate endpoint selection:
\begin{equation}
\boxed{
\begin{gathered}
\text{minimal endpoint class}\\
\Longrightarrow\ E_{\rm end}=\sigma V_3+\mu Q\\
\Longrightarrow\ \epsD(a)=\varepsilon_q+\varepsilon_w a^{-3},\quad \pD=-\varepsilon_q .
\end{gathered}
}
\end{equation}
Thus valuation closure constrains the background spectrum, the stated minimal endpoint class selects the vacuum-dust endpoint functional, and compact harmonics constrain the perturbative mode spectrum. They do not determine the dimensional amplitudes and do not establish statistical preference over \(\Lambda\)CDM. A predictive cosmological model additionally requires independent amplitude input, explicit perturbative realization, and direct likelihood-level comparison with current astronomical data.

\section{Observational Map, Superiority Criteria, and Falsification}
At background level, the minimal endpoint sector fixes the functional form of the residual sector within the stated endpoint class, but \(H(a)\) is determined only after the amplitudes and the visible sector have been specified. With radiation, baryons, the endpoint dust coefficient, and the vacuum coefficient, the minimal effective completion gives
\begin{equation}
H^2(a)=H_0^2\left[
\Omega_r a^{-4}
+\Omega_b a^{-3}
+\Omega_w a^{-3}
+\Omega_q
+\Omega_{K0} a^{-2}
\right],
\qquad
\Omega_{K0}=-\frac{c^2}{A_0^2H_0^2}<0
\end{equation}
for the closed convention used here. The density parameters satisfy
\begin{equation}
\Omega_r+\Omega_b+\Omega_w+\Omega_q+\Omega_{K0}=1,
\qquad
\Omega_m=\Omega_b+\Omega_w.
\end{equation}
Thus, at the minimal effective level, the observational background is not a new expansion law beyond curved \(\Lambda\)CDM with an effective cold component. The novelty is the full valuation classification, its finite dilation-spectrum closure, and the minimal endpoint classification of the residual functional form.

\subsection{Dilation-Closure Null Tests}
The differential characterization supplies sharp null tests once a residual density \(\epsD(a)\) has been reconstructed from the background under a declared visible-sector and curvature model. Define
\begin{align}
\mathcal R_{\rm val}[\epsD]
&:=(\Dil-1)\Dil(\Dil+1)(\Dil+2)(\Dil+3)\epsD,
\\
\mathcal R_{\rm bdry}[\epsD]
&:=\Dil(\Dil+1)(\Dil+2)(\Dil+3)\epsD,
\\
\mathcal R_{\rm end}[\epsD]
&:=\Dil(\Dil+3)\epsD.
\end{align}
The exact full valuation, boundary-only, and minimal endpoint classes satisfy respectively
\begin{equation}
\mathcal R_{\rm val}=0,
\qquad
\mathcal R_{\rm bdry}=0,
\qquad
\mathcal R_{\rm end}=0.
\end{equation}
Inside the full valuation class, the projected amplitudes must also satisfy
\begin{equation}
\frac{\dd\mathcal I_s}{\dd\ln a}=0,
\qquad
s\in\{1,0,-1,-2,-3\}.
\end{equation}
For the minimal endpoint sector the sharper conditions are
\begin{equation}
\mathcal I_1=\mathcal I_{-1}=\mathcal I_{-2}=0,
\qquad
\frac{\dd\mathcal I_0}{\dd\ln a}=0,
\qquad
\frac{\dd\mathcal I_{-3}}{\dd\ln a}=0.
\end{equation}

\begin{criterion}[Dilation-spectrum closure test]
A reconstructed residual history is consistent with a declared homothetic closure class only if the corresponding annihilator residual is statistically consistent with zero and the surviving projector amplitudes are consistent with constants over the tested interval, with uncertainties propagated through the full covariance of the reconstruction. Because high-order numerical differentiation is noise-sensitive, a practical implementation should infer these quantities from a covariance-aware basis fit, Gaussian-process reconstruction, or direct likelihood parameterization rather than from unregularized finite differences.
\end{criterion}

These null tests determine membership in the declared functional class; they do not by themselves identify a microscopic source or establish observational superiority.

\begin{nogotheorem}[Exact homogeneous-background degeneracy]
For the minimal endpoint completion, every homogeneous observable determined solely by \(H(z)\) and integrals of \(H^{-1}(z)\) is exactly degenerate with curved \(\Lambda\)CDM under
\begin{equation}
\Omega_m=\Omega_b+\Omega_w,
\qquad
\Omega_\Lambda=\Omega_q .
\end{equation}
Any statistical preference must therefore come from an additional prior, amplitude principle, compact spectral observable, perturbative observable, or non-minimal branch.
\end{nogotheorem}

\begin{proof}
The endpoint-sector background equation is
\begin{equation}
H^2(a)=H_0^2\left[
\Omega_r a^{-4}
+(\Omega_b+\Omega_w)a^{-3}
+\Omega_q
+\Omega_{K0}a^{-2}
\right].
\end{equation}
This is exactly the curved \(\Lambda\)CDM background form after the identifications
\begin{equation}
\Omega_m=\Omega_b+\Omega_w,
\qquad
\Omega_\Lambda=\Omega_q.
\end{equation}
Therefore any homogeneous background data set depending only on \(H(z)\), or on integrals of \(H^{-1}(z)\), cannot distinguish the minimal endpoint sector from curved \(\Lambda\)CDM after parameter relabelling. A statistical preference over curved \(\Lambda\)CDM must therefore come from information not contained in the homogeneous expansion law alone.
\end{proof}

\subsection{Background integrability}

Let
\begin{equation}
E^2(z):=\frac{H^2(z)}{H_0^2}.
\end{equation}
For fixed visible sector, fixed \(\Omega_{r0}\), fixed \(\Omega_{b0}\), fixed \(\Omega_{K0}\), and a separately conserved dark-energy branch with profile \(f_q(z)\), write
\begin{equation}
E^2(z)=
\Omega_{r0}(1+z)^4
+\Omega_{b0}(1+z)^3
+\Omega_{K0}(1+z)^2
+\Omega_{q0}f_q(z)
+\Omega_{w0}f_w(z).
\end{equation}
The reconstructed amplitude of the \(w\)-branch is
\begin{equation}
\widehat{\Omega}_{w0}(z)
=
\frac{
E^2(z)
-\Omega_{r0}(1+z)^4
-\Omega_{b0}(1+z)^3
-\Omega_{K0}(1+z)^2
-\Omega_{q0}f_q(z)
}{
f_w(z)
}.
\end{equation}

\begin{criterion}[Background integrability]
A single FRW background with one separately integrable \(w\)-branch requires
\begin{equation}
\frac{\dd}{\dd z}\widehat{\Omega}_{w0}(z)=0
\end{equation}
on the redshift interval on which the branch is defined and \(f_w(z)\neq0\). In a statistical implementation, the reconstructed values \(\widehat{\Omega}_{w0}(z_i)\) must be consistent with a single constant \(\Omega_{w0}\) under the full covariance matrix of the data.
\end{criterion}

For the endpoint dust branch,
\begin{equation}
f_w(z)=(1+z)^3.
\end{equation}
Failure of the constancy condition means that the proposed component is not one separately integrable dust-like branch of the assumed FRW background.

\subsection{Independent spectral-geometric invariant}

For a compact simply connected spatial section \(\Sthree\) with curvature radius \(R_c\), the standard scalar harmonic spectrum is \cite{HalliwellHawking1985,KodamaSasaki1984}
\begin{equation}
\operatorname{Spec}(-\nabla^2_{\Sthree})
=
\left\{
\frac{n(n+2)}{R_c^2}
\right\}_{n=0}^{\infty},
\qquad
 d_n=(n+1)^2,
\end{equation}
where \(d_n\) is the scalar-mode degeneracy. With the closed-sign convention \(\Omega_{K0}<0\),
\begin{equation}
R_c=\frac{c}{H_0\sqrt{-\Omega_{K0}}}.
\end{equation}

A smooth \(w(z)\)-parameterization may change the inferred distance-redshift relation, growth history, or preferred curvature radius. It does not replace the discrete compact-\(\Sthree\) eigenvalue structure by the continuous flat-space spectrum.

\begin{criterion}[Spectral-geometric independence]
A compact-\(\Sthree\) branch is independently supported only if the inferred curvature scale and perturbation analysis are compatible with the discrete scalar, vector, and tensor harmonic spectra of \(\Sthree\). In the simply connected case this concerns curvature scale, mode discreteness, and large-angle correlations. Matched circles or multiple imaging require an additional nontrivial quotient \(\Sthree/\Gamma\) and are not consequences of simply connected \(\Sthree\) alone.
\end{criterion}

Cosmic-topology searches and matched-circle tests belong to this quotient-space sector rather than to simply connected \(\Sthree\) alone \cite{CornishSpergelStarkman1998,LuminetEtAl2003,Planck2016Topology}.

Absence of a detected spectral imprint does not by itself disprove compactness if the curvature radius is too large for the data to resolve. It does, however, remove independent confirmation of the compact-\(\Sthree\) branch unless the completed model predicts no observable spectral effect in the tested range.

\begin{remark}[What the compact spectrum can and cannot distinguish]
The compact \(\Sthree\) spectrum is a geometric and topological discriminator. It can distinguish a compact curved spatial section from a flat or non-compact parametrization, provided the curvature scale is observationally accessible. It does not by itself distinguish the endpoint matter law from a closed curved \(\Lambda\)CDM model placed on the same \(\Sthree\) geometry. Endpoint-sector distinguishability requires additional information: amplitude priors, perturbative observables, spectral weights, allowed initial spectra, non-minimal valuation branches, or penalized statistical preference against competitors with the same spatial topology included.
\end{remark}

\subsection{Robust statistical separation}

Let \(X\) run over the nearest competing models,
\begin{equation}
X\in
\left\{
\Lambda{\rm CDM},\;
\Lambda{\rm CDM}+\Omega_{K0},\;
w_0w_a{\rm CDM},\;
w_0w_a{\rm CDM}+\Omega_{K0}
\right\}.
\end{equation}
For the same data vector, likelihood construction, nuisance treatment, and parameter priors, define
\begin{align}
\Delta{\rm AIC}_X
&=
{\rm AIC}_{\rm model}-{\rm AIC}_X,
\\
\Delta{\rm BIC}_X
&=
{\rm BIC}_{\rm model}-{\rm BIC}_X,
\\
\Delta\ln Z_X
&=
\ln Z_{\rm model}-\ln Z_X .
\end{align}

\begin{criterion}[Robust statistical branch]
The proposed branch is statistically autonomous against the nearest competitors only if
\begin{equation}
\Delta{\rm AIC}_X\le 0,
\qquad
\Delta{\rm BIC}_X\le 0,
\qquad
\Delta\ln Z_X\ge 0
\end{equation}
for every listed competitor \(X\), using the same data choices and parameter accounting. The conclusion must remain stable under replacement of supernova samples and under the addition of \(f\sigma_8\), weak-lensing, and CMB-lensing data.
\end{criterion}

Planck 2018 combined with BAO found spatial curvature consistent with flatness, \(\Omega_{K0}=0.001\pm0.002\), so closed curvature is not supported by that combined constraint. DESI DR2 provides current BAO distance constraints from the first three years of DESI observations and is a proper reference point for testing \(\Lambda\)CDM and dark-energy extensions, but it does not by itself establish a closed \(\Sthree\) endpoint cosmology \cite{Planck2018Parameters,DESI2025DR2BAO}. Therefore any claim of superiority must specify the data set, likelihood, priors, nuisance model, and competitor set.

\subsection{Conditional superiority criterion}

\begin{criterion}[Conditional observational superiority]
Within the declared FRW model class, the completed branch is observationally superior to the listed competitors only if all three conditions hold:
\begin{enumerate}
\item the reconstructed branch amplitude is integrable,
\begin{equation}
\frac{\dd}{\dd z}\widehat{\Omega}_{w0}(z)=0
\end{equation}
within the joint covariance of the data;
\item the compact-\(\Sthree\) sector has independent spectral-geometric support, or the model supplies a justified reason why the compact spectral signature is beyond the tested observational range;
\item the penalized model comparison satisfies
\begin{equation}
\Delta{\rm AIC}_X\le0,
\qquad
\Delta{\rm BIC}_X\le0,
\qquad
\Delta\ln Z_X\ge0
\end{equation}
for every nearest competitor \(X\).
\end{enumerate}
If any condition fails, observational superiority has not been established.
\end{criterion}

\begin{proof}
The first condition excludes a redshift-dependent phenomenological remainder masquerading as a single FRW component. The second condition supplies an invariant compact-geometric datum independent of a smooth \(w(z)\)-fit. The third condition prevents an apparent improvement in likelihood from being only a penalty-unaware increase in model flexibility. Together these conditions separate a physical branch from a phenomenological fit inside the declared FRW class.
\end{proof}

\subsection{Falsification logic}

The construction is falsified or restricted at the corresponding level if at least one of the following holds:
\begin{enumerate}
\item the endpoint-selected or independently selected \(\Efour(A,\xi)\) fails to reproduce the reconstructed residual in the matching condition;
\item the internal equations are incompatible with the exchange law;
\item the dust slot necessarily has unacceptable \(c_s^2\), \(\pi_{\mu\nu}\), or perturbation growth;
\item the compact radius required by an independently completed model is excluded by curvature, spectral, or topology tests;
\item the gluing conditions of the compatible boundary-data condition are incompatible with the proposed \(\Sfour\) functional;
\item no covariant completion exists, or the completion violates the stated symmetries;
\item the reconstructed branch amplitude fails the integrability condition;
\item the compact-\(\Sthree\) branch has no independent spectral-geometric support in the tested data regime;
\item the completed branch fails robust penalized comparison against the nearest competitors.
\end{enumerate}

Thus background integrability, spectral-geometric discreteness, and robust statistical performance are necessary filters for turning a geometric branch into an observationally superior physical cosmology.


\section{Status of Results and Remaining Program}
\begin{center}
\begin{tabular}{@{}p{0.20\linewidth}p{0.50\linewidth}p{0.22\linewidth}@{}}
\toprule
Object & Precise result & Status \\
\midrule
\(\Bfour/\Sthree\) geometry & \(V_4/V_3=A/4\), including fixed quasilocal cones & proved \\
Direct map & \(\Efour\mapsto(\epsD,\pD)\) and background conservation & proved \\
Inverse representation & \(\epsD\mapsto\rhoFour=4\epsD/A\) in the homogeneous class & proved \\
Equation of state & General \(w_D\) formula, power-law classes, and energy conditions & proved \\
Dimensional structure & \(w_d=(\beta-d-1)/d\) & proved \\
Superposition & Exact addition of background components & proved \\
Full valuation spectrum & \(\epsD=\gamma A/4+\alpha_0+3\alpha_1/A+3\alpha_2/A^2+\alpha_3/A^3\) & proved under stated valuation assumptions \\
Dilation annihilator & \((\Dil-1)\Dil(\Dil+1)(\Dil+2)(\Dil+3)\epsD=0\) iff round valuation profile image & proved \\
Endpoint density obstruction & \(\Dil(\Dil+3)\epsD=0\) iff vacuum-dust density profile; endpoint lift is separate & proved \\
Pressure kernel & \(p=0\iff E=\mathrm{const}\Rightarrow\varepsilon\propto a^{-3}\) & proved \\
Boundary-only valuation restriction & \(\gamma=0\) & conditional within boundary-only valuation class \\
Minimal endpoint classification & Inside the stated class, \(E_{\rm end}=\sigma V_3+\mu Q\Rightarrow \epsD=\varepsilon_q+\varepsilon_w a^{-3}\) & conditional theorem within the minimal endpoint class \\
Vacuum and dust & \(\epsD=\varepsilon_q+\varepsilon_w a^{-3}\), \(\pD=-\varepsilon_q\) & selected by minimal endpoint classification \\
Amplitude determination & Numerical values of \(\sigma\) and \(\mu\) & not fixed by topology \\
Minimal covariant completion & Vacuum plus conserved dust current, with \(\mu Q=2\pi^2mn_0A_0^3\) & constructed \\
Dust perturbations & \(c_s^2=0\), \(\pi_{\mu\nu}=0\), geodesic flow & satisfied by minimal completion \\
\(T^{\rm fill}=T^{\rm rec}\) & Independent reproduction of the residual & model condition \\
\(\Sfour\) compensation & Boundary-form cancellation; Chern--Simons phase does not fix \(\sigma\) & conditionally proved plus no-go \\
Hopf and causality & Candidate internal variables are defined & hypotheses \\
Radius & Topological no-go theorem & proved \\
Gauss--Kronecker constant slot & Gauss-map realization of the \(A^0\) valuation slot, not an independent \(K\)-action in the minimal endpoint class & clarified \\
Topological charge option & \(N[U]\in\mathbb{Z}\); dust tensor requires action matching & conditional completion \\
Observational autonomy & Integrability, spectra, perturbations, penalized statistics, and calibration & remaining program \\
\bottomrule
\end{tabular}
\end{center}

The remaining nontrivial tasks are:
\begin{enumerate}
\item justify, test, or replace the boundary-only valuation restriction and the assumptions defining the minimal endpoint class;
\item determine or calibrate the dimensional amplitudes \(\sigma\) and \(\mu\);
\item construct any non-minimal four-dimensional diffeomorphism-invariant completion and verify that it reduces to the endpoint-selected sector;
\item identify additional internal fields \(\xi^I\), if introduced, together with their gauge structure and second-theorem Noether identities;
\item construct a concrete \(\Sfour\) compensation functional with its dimensionful couplings if amplitude selection by compensation is claimed, since Chern--Simons gluing alone fixes only topological phases;
\item if a topological charge completion is used, supply the field \(U:\M\to SU(2)\), its action, its conserved charge, and the dimensional mass scale needed to match \(\mu\);
\item perform the observational analysis with a dated data set, including curvature scale, harmonic spectra, structure growth, lensing, CMB constraints, background integrability, and penalized comparison against the nearest competing cosmologies, with competitors of the same spatial topology included when relevant.
\end{enumerate}





\section{Closure of the Program at the Present Level}

The \(\Bfour/\Sthree\) framework has four logically distinct layers.

The first layer is the universal homogeneous background representation. Every differentiable \(\epsD(A)\) has the unique representative
\begin{equation}
\rhoFour(A)=\frac{4}{A}\epsD(A)
\end{equation}
in the adopted homothetic filling class. This layer is mathematically complete but not predictive by itself.

The second layer is the full homothetic valuation classification. Under the stated Euclidean-invariant valuation assumptions for the auxiliary convex filling in \(\mathbb{R}^4\), the available powers are
\begin{equation}
\epsD(A)=
\frac{\gamma A}{4}
+\alpha_0+\frac{3\alpha_1}{A}
+\frac{3\alpha_2}{A^2}
+\frac{\alpha_3}{A^3},
\end{equation}
with branch equations of state
\begin{equation}
w=-\frac{4}{3},\ -1,\ -\frac{2}{3},\ -\frac{1}{3},\ 0.
\end{equation}
This is the theorem-level classification result beyond the universal map. Its invariant normal form is
\begin{equation}
(\Dil-1)\Dil(\Dil+1)(\Dil+2)(\Dil+3)\epsD=0,
\qquad
\Dil:=A\frac{\dd}{\dd A}.
\end{equation}
Thus the valuation class is a finite spectral subspace of the homothetic dilation operator rather than merely a list of fitted powers. The spectral projectors \(\mathsf P_s\) extract invariant amplitudes \(\mathcal I_s\); in particular the bulk obstruction is \(\mathcal I_{\rm bulk}=\mathcal I_1\).

The third layer is endpoint selection inside a specified minimal boundary class. Independently of its topological realization, the endpoint density has the sharp differential characterization
\begin{equation}
\Dil(\Dil+3)\epsD=0
\quad\Longleftrightarrow\quad
\epsD=\varepsilon_q+\varepsilon_w a^{-3}.
\end{equation}
The surviving amplitudes are the dilation invariants \(\varepsilon_q=\mathcal I_0\) and \(\varepsilon_w=\mathcal I_{-3}\). Moreover, the dust term is the image of the constant-energy kernel of the work-pressure map under division by \(V_3\): \(p=0\iff E=\mathrm{const}\Rightarrow\varepsilon\propto a^{-3}\).

The endpoint form is not a theorem from topology alone. It is a conditionally unique theorem inside the minimal endpoint class: local \({\rm Diff}(\Sthree)\)-invariant boundary functionals with no local propagating endpoint degrees of freedom, no gravitational curvature duplication, no fixed external geometric structures, no background connections, no nonlocal spectral invariants, no additional geometric or dynamical length scales in the functional dependence, additivity, and internal-field dependence only through an oriented topological charge. Inside that class,
\begin{equation}
E_{\rm end}(A,Q)=\sigma V_3(A)+\mu Q,
\end{equation}
and hence
\begin{equation}
\epsD(a)=\varepsilon_q+\varepsilon_w a^{-3},
\qquad
\pD=-\varepsilon_q,
\end{equation}
with
\begin{equation}
\varepsilon_q=\sigma,
\qquad
\varepsilon_w=\frac{\mu Q}{2\pi^2A_0^3}.
\end{equation}
The boundary-only valuation restriction \(\gamma=0\) remains an independent geometric filter, but it is not the sole mechanism by which the endpoint background form is isolated.

A future topological-charge completion may identify \(Q\) with a winding number of a map \(U:\Sthree\to SU(2)\), but the dimensional mass-per-charge scale remains independent. Topology supplies the integer charge; it does not supply \(\mu\). If the covariant current is built from \(U:\M\to SU(2)\), it becomes an effective cold-dust tensor only after the explicit pressureless action matching given above.

The fourth layer is observational autonomy. The minimal endpoint completion is vacuum energy plus a conserved dust current. It realizes the selected dust branch with
\begin{equation}
c_s^2=0,
\qquad
\pi_{\mu\nu}=0,
\qquad
u^\mu\nabla_\mu u^\nu=0.
\end{equation}
This proves existence of a perturbatively cold completion, not uniqueness of the microscopic realization.

At the homogeneous observational level, the minimal endpoint sector is exactly curved \(\Lambda\)CDM after the identification
\begin{equation}
\Omega_m=\Omega_b+\Omega_w,
\qquad
\Omega_\Lambda=\Omega_q.
\end{equation}
Therefore the minimal endpoint sector cannot claim background-only superiority over curved \(\Lambda\)CDM. Any genuine superiority must come from at least one additional structure: compact-\(\Sthree\) spectra, topology or curvature priors, non-minimal valuation branches, an amplitude principle, or perturbative observables.

Consequently, the strongest defensible claim is:
\begin{quote}
The work classifies the full homothetic valuation spectrum of the \(\Bfour/\Sthree\) filling and proves that, inside the stated minimal endpoint class, the endpoint functional has the conditional form \(E_{\rm end}=\sigma V_3+\mu Q\). Its minimal four-dimensional effective realization is curved \(\Lambda\)CDM with cold dust perturbations. Observational superiority is not assumed; it is promoted to a falsifiable data-dependent claim requiring background integrability, independent compact-\(\Sthree\) spectral support, and robust penalized statistical performance against the nearest competing cosmologies.
\end{quote}

The construction should not be read as a derivation of the numerical amplitudes, a unique microscopic matter model, or a complete dynamics for \(A(\tau)\). Chern--Simons gluing can constrain oriented boundary phases but cannot determine the vacuum density \(\sigma\) without an independent dimensionful input. Numerical amplitudes therefore require an additional dimensional principle, coupling, quantum input, or observational calibration.

\section{Conclusion}

A rigorous background \(\Bfour/\Sthree\) representation map has been constructed:
\begin{equation}
\Efour(A,\xi)\longmapsto(\epsD,\pD,\QD).
\end{equation}
The radial derivative of the filling energy has been assigned its correct status as a conjugate response; pressure has been derived from work; conservation has been obtained as an identity; the direct map has been complemented by the inverse representation theorem; the power-law relation has been generalized to a local \(w_D(A)\); and both the dimensional theorem and linear superposition have been proved.

The central theorem-level strengthening is the full homothetic valuation spectrum:
\begin{equation}
\boxed{
\text{full valuation class}
\quad\Rightarrow\quad
\epsD(A)=
\frac{\gamma A}{4}
+\alpha_0+\frac{3\alpha_1}{A}
+\frac{3\alpha_2}{A^2}
+\frac{\alpha_3}{A^3}
}
\end{equation}
with branches
\begin{equation}
w=-\frac{4}{3},\ -1,\ -\frac{2}{3},\ -\frac{1}{3},\ 0.
\end{equation}
The theorem-level invariant core is not the existence of the algebraic map
\begin{equation}
\Efour\longmapsto(\epsD,\pD),
\end{equation}
since that map admits an inverse for every differentiable homogeneous density in the adopted class. The predictive closure is the finite dilation spectrum selected by the round homothetic valuation image:
\begin{equation}
\operatorname{Spec}_{\rm val}(\Dil)=\{1,0,-1,-2,-3\},
\qquad
\Dil:=A\frac{\dd}{\dd A}.
\end{equation}
Equivalently,
\begin{equation}
\boxed{
(\Dil-1)\Dil(\Dil+1)(\Dil+2)(\Dil+3)\epsD=0.
}
\end{equation}
This gives a sharp membership test for the round homothetic profile image of the full valuation class and invariant projectors \(\mathcal I_s\) for its five amplitudes. The boundary-only valuation restriction is equivalently the obstruction condition \(\mathcal I_{\rm bulk}=\mathcal I_1=0\). The minimal endpoint density is the two-eigenvalue spectral subspace
\begin{equation}
\boxed{
\ker[\Dil(\Dil+3)]
=\operatorname{span}\{a^0,a^{-3}\}.
}
\end{equation}
Its invariant amplitudes are \(\varepsilon_q=\mathcal I_0\) and \(\varepsilon_w=\mathcal I_{-3}\). The endpoint theorem supplies the separate conditional functional lift from this density subspace to \(E_{\rm end}=\sigma V_3+\mu Q\) inside the stated minimal endpoint class. The pressureless branch has the additional structural characterization
\begin{equation}
p=0\iff E=\mathrm{const},
\end{equation}
so the dust scaling is the image of the constant-energy kernel of the work-pressure map under division by \(V_3\).

The boundary-only valuation restriction gives \(\gamma=0\) as an independent geometric filter. Inside the stated minimal endpoint class one obtains
\begin{equation}
\boxed{
E_{\rm end}(A,Q)=\sigma V_3(A)+\mu Q
\quad\Rightarrow\quad
\epsD(a)=\varepsilon_q+\varepsilon_w a^{-3},
\qquad
\pD=-\varepsilon_q
}
\end{equation}
with
\begin{equation}
\varepsilon_q=\sigma,
\qquad
\varepsilon_w=\frac{\mu Q}{2\pi^2A_0^3}.
\end{equation}
The endpoint form is not a theorem from topology alone. It is a conditionally unique theorem inside the minimal endpoint class: local \({\rm Diff}(\Sthree)\)-invariant boundary functionals with no local propagating endpoint degrees of freedom, no gravitational curvature duplication, no fixed external geometric structures, no background connections, no nonlocal spectral invariants, no additional geometric or dynamical length scales in the functional dependence, additivity, and internal-field dependence only through an oriented topological charge.

The amplitude no-go theorem fixes the exact scope of this predictivity: pure topology, the volume ratio, the valuation basis, and the endpoint invariants classify or select powers of \(A\), not the numerical energy scales. A minimal covariant effective completion by vacuum energy plus a conserved dust current realizes the selected dust slot with
\begin{equation}
c_s^2=0,
\qquad
\pi_{\mu\nu}=0,
\qquad
u^\mu\nabla_\mu u^\nu=0,
\end{equation}
and matching to the endpoint constant gives
\begin{equation}
\mu Q=2\pi^2mn_0A_0^3.
\end{equation}
This establishes a cold effective perturbative completion of the selected sector. If a later topological completion is added, it may quantize a conserved number \(Q\), but the dimensional amplitude still requires an independent mass scale. The Gauss--Kronecker endpoint term is likewise a Gauss-map degree invariant in the round convex case, not an Euler characteristic of \(\Sthree\). Chern--Simons gluing supplies oriented phase consistency, not the vacuum amplitude \(\sigma\) by itself.

At the minimal effective observational level the model is curved \(\Lambda\)CDM with an effective cold component. Therefore it is not a new background expansion law and cannot claim background-only superiority over curved \(\Lambda\)CDM. Its possible distinguishing content lies in the full valuation classification, the minimal endpoint classification, compact \(\Sthree\) spectra, topology or curvature priors, non-minimal branches, perturbative observables, or an amplitude principle.

The final falsifiable criterion is the following:
\begin{equation}
\begin{array}{l}
\text{background integrability}\\[1mm]
+\ \text{independent compact-}\Sthree\text{ spectral-geometric support}\\[1mm]
+\ \text{robust penalized statistical gain against nearest competitors}
\end{array}
\end{equation}
is required before observational superiority can be asserted. Failure of any one of these conditions means that the construction remains a geometric classification and effective completion, not an empirically superior cosmology.

The final status of the construction is therefore deliberately limited. It is invariantly admissible, compact-compatible, background-reconstructive, conditionally endpoint-selective, and perturbatively closed for the minimal vacuum-plus-dust effective realization. It is not, by itself, an amplitude derivation, a microscopic source theory, a transfer-function calculation, or an observationally autonomous cosmology. In particular,
\begin{equation}
\begin{array}{c}
\text{residual}\neq\text{source}\neq\text{amplitude}\\[1mm]
\neq\text{transfer}\neq\text{likelihood}.
\end{array}
\end{equation}
The strongest theorem-level content is the chain
\begin{equation}
\begin{array}{c}
\Bfour_A/\Sthree_A\ \text{homothetic representation}\\[1mm]
\Longrightarrow\ \text{finite valuation spectrum}\\[1mm]
\Longrightarrow\ \text{conditional minimal endpoint form}.
\end{array}
\end{equation}
With
\begin{equation}
E_{\rm end}(A,Q)=\sigma V_3(A)+\mu Q,
\qquad
\epsD(a)=\varepsilon_q+\varepsilon_w a^{-3},
\qquad
\pD=-\varepsilon_q .
\end{equation}
Thus the framework is benchmark-useful and falsifiable, but amplitude selection, microscopic origin, transfer functions, likelihood construction, and observational covariance remain independent completion problems.

This does not identify \(\Bfour_A\) with a physical bulk, the effective residual sector with a unique microscopic matter model, or the limiting behaviour of \(A\) with an ontological endpoint of the carrier. The remaining physical inputs are the amplitude scale, possible non-minimal completions outside the minimal endpoint class, and observational selection.

\section*{Data and Code Availability}
This theoretical work generated no new data sets and used no numerical code to obtain its analytic results.

\begin{thebibliography}{99}

\bibitem{Noether1918}
E. Noether, ``Invariante Variationsprobleme,'' \emph{Nachrichten von der Gesellschaft der Wissenschaften zu G\"ottingen, Mathematisch-Physikalische Klasse} (1918), 235--257.

\bibitem{LeeWald1990}
J. Lee and R. M. Wald, ``Local symmetries and constraints,'' \emph{Journal of Mathematical Physics} \textbf{31} (1990), 725--743, \href{https://doi.org/10.1063/1.528801}{doi:10.1063/1.528801}.

\bibitem{BrownYork1993}
J. D. Brown and J. W. York, ``Quasilocal energy and conserved charges derived from the gravitational action,'' \emph{Physical Review D} \textbf{47} (1993), 1407--1419, \href{https://doi.org/10.1103/PhysRevD.47.1407}{doi:10.1103/PhysRevD.47.1407}.

\bibitem{IyerWald1994}
V. Iyer and R. M. Wald, ``Some properties of Noether charge and a proposal for dynamical black hole entropy,'' \emph{Physical Review D} \textbf{50} (1994), 846--864, \href{https://doi.org/10.1103/PhysRevD.50.846}{doi:10.1103/PhysRevD.50.846}.

\bibitem{BarnichBrandt2002}
G. Barnich and F. Brandt, ``Covariant theory of asymptotic symmetries, conservation laws and central charges,'' \emph{Nuclear Physics B} \textbf{633} (2002), 3--82, \href{https://doi.org/10.1016/S0550-3213(02)00251-1}{doi:10.1016/S0550-3213(02)00251-1}.

\bibitem{GibbonsHawking1977}
G. W. Gibbons and S. W. Hawking, ``Action integrals and partition functions in quantum gravity,'' \emph{Physical Review D} \textbf{15} (1977), 2752--2756, \href{https://doi.org/10.1103/PhysRevD.15.2752}{doi:10.1103/PhysRevD.15.2752}.

\bibitem{DeWitt1967}
B. S. DeWitt, ``Quantum theory of gravity. I. The canonical theory,'' \emph{Physical Review} \textbf{160} (1967), 1113--1148, \href{https://doi.org/10.1103/PhysRev.160.1113}{doi:10.1103/PhysRev.160.1113}.

\bibitem{HawkingEllis1973}
S. W. Hawking and G. F. R. Ellis, \emph{The Large Scale Structure of Space-Time}, Cambridge University Press, Cambridge, 1973.

\bibitem{MTW1973}
C. W. Misner, K. S. Thorne, and J. A. Wheeler, \emph{Gravitation}, W. H. Freeman, San Francisco, 1973.

\bibitem{Weinberg2008}
S. Weinberg, \emph{Cosmology}, Oxford University Press, Oxford, 2008.

\bibitem{KodamaSasaki1984}
H. Kodama and M. Sasaki, ``Cosmological perturbation theory,'' \emph{Progress of Theoretical Physics Supplement} \textbf{78} (1984), 1--166, \href{https://doi.org/10.1143/PTPS.78.1}{doi:10.1143/PTPS.78.1}.

\bibitem{HalliwellHawking1985}
J. J. Halliwell and S. W. Hawking, ``Origin of structure in the Universe,'' \emph{Physical Review D} \textbf{31} (1985), 1777--1791, \href{https://doi.org/10.1103/PhysRevD.31.1777}{doi:10.1103/PhysRevD.31.1777}.

\bibitem{MukhanovFeldmanBrandenberger1992}
V. F. Mukhanov, H. A. Feldman, and R. H. Brandenberger, ``Theory of cosmological perturbations,'' \emph{Physics Reports} \textbf{215} (1993), 203--333, \href{https://doi.org/10.1016/0370-1573(92)90044-Z}{doi:10.1016/0370-1573(92)90044-Z}.

\bibitem{MalikWands2009}
K. A. Malik and D. Wands, ``Cosmological perturbations,'' \emph{Physics Reports} \textbf{475} (2009), 1--51, \href{https://doi.org/10.1016/j.physrep.2009.03.001}{doi:10.1016/j.physrep.2009.03.001}.

\bibitem{Hopf1931}
H. Hopf, ``\"Uber die Abbildungen der dreidimensionalen Sph\"are auf die Kugelfl\"ache,'' \emph{Mathematische Annalen} \textbf{104} (1931), 637--665, \href{https://doi.org/10.1007/BF01457962}{doi:10.1007/BF01457962}.

\bibitem{CornishSpergelStarkman1998}
N. J. Cornish, D. N. Spergel, and G. D. Starkman, ``Circles in the sky: finding topology with the microwave background radiation,'' \emph{Classical and Quantum Gravity} \textbf{15} (1998), 2657--2670, \href{https://doi.org/10.1088/0264-9381/15/9/005}{doi:10.1088/0264-9381/15/9/005}.

\bibitem{LuminetEtAl2003}
J.-P. Luminet, J. R. Weeks, A. Riazuelo, R. Lehoucq, and J.-P. Uzan, ``Dodecahedral space topology as an explanation for weak wide-angle temperature correlations in the cosmic microwave background,'' \emph{Nature} \textbf{425} (2003), 593--595, \href{https://doi.org/10.1038/nature01944}{doi:10.1038/nature01944}.

\bibitem{Hadwiger1957}
H. Hadwiger, \emph{Vorlesungen \"uber Inhalt, Oberfl\"ache und Isoperimetrie}, Springer, Berlin, 1957.

\bibitem{Schneider2014}
R. Schneider, \emph{Convex Bodies: The Brunn-Minkowski Theory}, 2nd expanded ed., Cambridge University Press, Cambridge, 2014.

\bibitem{Brown1993Fluid}
J. D. Brown, ``Action functionals for relativistic perfect fluids,'' \emph{Classical and Quantum Gravity} \textbf{10} (1993), 1579--1606, \href{https://doi.org/10.1088/0264-9381/10/8/017}{doi:10.1088/0264-9381/10/8/017}.

\bibitem{ChernSimons1974}
S.-S. Chern and J. Simons, ``Characteristic forms and geometric invariants,'' \emph{Annals of Mathematics} \textbf{99} (1974), 48--69, \href{https://doi.org/10.2307/1971013}{doi:10.2307/1971013}.

\bibitem{Skyrme1961}
T. H. R. Skyrme, ``A non-linear field theory,'' \emph{Proceedings of the Royal Society of London A} \textbf{260} (1961), 127--138, \href{https://doi.org/10.1098/rspa.1961.0018}{doi:10.1098/rspa.1961.0018}.

\bibitem{Witten1983CurrentAlgebra}
E. Witten, ``Global aspects of current algebra,'' \emph{Nuclear Physics B} \textbf{223} (1983), 422--432, \href{https://doi.org/10.1016/0550-3213(83)90063-9}{doi:10.1016/0550-3213(83)90063-9}.

\bibitem{Planck2018Parameters}
Planck Collaboration, ``Planck 2018 results. VI. Cosmological parameters,'' \emph{Astronomy \& Astrophysics} \textbf{641} (2020), A6, \href{https://doi.org/10.1051/0004-6361/201833910}{doi:10.1051/0004-6361/201833910}.

\bibitem{DESI2025DR2BAO}
DESI Collaboration, M. Abdul Karim et al., ``DESI DR2 Results II: Measurements of Baryon Acoustic Oscillations and Cosmological Constraints,'' arXiv:2503.14738 [astro-ph.CO].

\bibitem{Planck2016Topology}
Planck Collaboration, ``Planck 2015 results. XVIII. Background geometry and topology of the Universe,'' \emph{Astronomy \& Astrophysics} \textbf{594} (2016), A18, \href{https://doi.org/10.1051/0004-6361/201525829}{doi:10.1051/0004-6361/201525829}.

\end{thebibliography}

\end{document}
